\documentclass[11pt,a4paper,oneside]{book}
\usepackage[T1]{fontenc}
\usepackage[utf8]{inputenc}
\usepackage{lmodern,microtype}
\usepackage[margin=28mm]{geometry}
\usepackage{amsmath,amssymb,amsthm,mathtools}
\usepackage{enumitem}
\usepackage{xcolor}
\usepackage{xurl}
\usepackage[hidelinks]{hyperref}
\usepackage{bookmark}
\usepackage{fancyvrb}
\allowdisplaybreaks
\setlength{\emergencystretch}{3em}
\newtheorem{theorem}{Theorem}[chapter]
\newtheorem{lemma}[theorem]{Lemma}
\newtheorem{proposition}[theorem]{Proposition}
\newtheorem{corollary}[theorem]{Corollary}
\theoremstyle{definition}
\newtheorem{definition}[theorem]{Definition}
\newtheorem{remark}[theorem]{Remark}
\newtheorem{analyticinput}[theorem]{Established analytic input}
\newcommand{\R}{\mathbb R}
\newcommand{\C}{\mathbb C}
\newcommand{\N}{\mathbb N}
\newcommand{\Z}{\mathbb Z}
\newcommand{\eps}{\varepsilon}
\newcommand{\ii}{\mathrm i}
\newcommand{\e}{\mathrm e}
\newcommand{\one}{\mathbf 1}
\newcommand{\abs}[1]{\left|#1\right|}
\newcommand{\norm}[1]{\left\|#1\right\|}
\newcommand{\tr}{\operatorname{tr}}
\newcommand{\supp}{\operatorname{supp}}
\newcommand{\Rea}{\operatorname{Re}}
\newcommand{\Ima}{\operatorname{Im}}
\newcommand{\lean}[1]{\nolinkurl{#1}}
\newcommand{\source}[1]{\par\nobreak\smallskip\nobreak\noindent{\footnotesize\textit{Lean correspondence:}\ \lean{#1}.}\par\smallskip}
\newcommand{\Hhat}{\widehat}
\title{Zero-density estimates from an improved large-values bound\\[6pt]\large A source-faithful mathematical blueprint of the registered Lean development}
\author{Mathematics formalized by Mark Lewko\\\normalsize Informal reconstruction for study; not an author-approved paper}
\date{Source: Palomar entry PALOMAR-2026-10-11-000010, version 1}
\begin{document}
\pagestyle{plain}
\frontmatter
\hypersetup{pageanchor=false}
\maketitle
\hypersetup{pageanchor=true}
\chapter*{Scope, provenance, and reading conventions}
\addcontentsline{toc}{chapter}{Scope, provenance, and reading conventions}
This document reconstructs the mathematical argument represented by the four
endpoint theorems in \lean{lean/Solution.lean} of
\lean{mrklwk/zeta-density-39-over-17-lean}, at the exact commit
\begin{center}\lean{ee2a8a32be49a4f73c7386e9dffdea24d7e3d8c0}.\end{center}
The immutable record is
\url{https://data.palomar-registry.org/entries/PALOMAR-2026-10-11-000010-v1.json};
the requested entry is
\url{https://palomar-registry.org/entry?id=PALOMAR-2026-10-11-000010&version=1}.
The dependency on mathlib is pinned to
\lean{065356127b1dc0016f66b7283ce0ce2c4055aa55}.
All references to a source file mean this snapshot, not a changing default branch.

\paragraph{What is reconstructed.}
The text gives full proofs of the endpoint composition, the middle-range scale
optimization and absorption, the passage from truncated moments to separated
values, the two new density windows, multiplicity restoration, dyadic summation,
and uniformization. It also spells out the operators, Fourier masks, affine
averages, detector families, and imported analytic interfaces needed by those
proofs. The substantial previously established analytic machinery is explicitly
identified as such: an invocation of an established input below is not a new
assumption in an endpoint theorem. In particular, a proof in this blueprint is
understood relative to the precisely stated, already proved inputs that it cites.
The native fourth-moment, point-mean/stationary-phase, classical reflection, and
Fourier-affine foundations are not represented as if they followed from elementary
algebra. Their actual role and source interfaces are recorded separately.

\paragraph{Boundary of the transcription.}
This is a mathematical blueprint, not a literal translation of every auxiliary
lemma in all 824 locally imported modules. There are supporting library proofs
behind the established inputs. These are not re-proved down to the foundations
of measure theory, special functions, or the native stationary-source ports here.
The accompanying complete pinned source archive preserves those proofs. Source
correspondences make this boundary explicit, rather than concealing it with
phrases such as ``standard estimates.'' No independent verification of a fresh
full Lean build is claimed by this reconstruction.

\paragraph{Fidelity rules.}
All real parameters in the endpoint statements are retained as real parameters.
The interval of source integers is exactly $(N,2N]$, implemented with natural
floors. Both signs of zero ordinates and both closed height endpoints are
retained. Zero locations and zero multiplicities are never identified. A fixed
constant is selected before every variable over which the corresponding result
is uniform. The rightmost density branch starts \emph{strictly} above $3/4$.
We do not assert a globally uniform constant for the exact discontinuous curve.
The historical name \lean{CandidateMiddleRangeEstimate} denotes a proposition
that is proved; conversely, the separately defined
\lean{ShortMobiusFourthMoment} is not silently used as a theorem.

\paragraph{Logical and analytic conventions.}
Integrals on $\R$ use Lebesgue measure unless another measure is displayed.
The Fourier transform is
\[
 \Hhat f(\xi)=\int_{\R}\e^{-2\pi\ii\xi t}f(t)\,dt,
 \qquad e(x)=\e^{2\pi\ii x},\qquad
 \lambda(x)=\frac{\log x}{2\pi}\quad(x>0).
\]
Its Plancherel constant is one. Complex inner products are conjugate-linear in
the first argument. A real power of a positive real base means
$x^a=\exp(a\log x)$. Every application below of monotonicity in an exponent has
base at least one. ``Eventually in $H$'' means that one height threshold exists
before all subsequently quantified variables. A notation $\ll_{\eps}$ abbreviates
an existential positive constant independent of all other displayed varying
objects, not an effective numerical value.

\paragraph{Proof automation.}
Source uses of \lean{ring}, \lean{field_simp}, \lean{norm_num},
\lean{linarith}, \lean{nlinarith}, and \lean{gcongr} are expanded as algebra,
positive-denominator comparisons, or monotonicity. They are not additional
analytic inputs. Classical finite-set selection and Hilbert-basis selection
account for uses of choice. The registered axiom allowance is
\lean{propext}, \lean{Classical.choice}, and \lean{Quot.sound}.
The holes in the independent \lean{Challenge.lean} specification are not imported
by \lean{Solution.lean} and are not premises of the proof.
\tableofcontents
\mainmatter

\chapter{The objects counted and the four endpoint statements}
\section{Actual zeros and analytic multiplicity}
\begin{definition}
Let $\zeta$ be mathlib's Riemann zeta function. A nontrivial zero is a complex
number $\rho$ for which
\[
 \zeta(\rho)=0,\qquad 0<\Rea\rho<1.
\]
Define $m(\rho)$ to be the natural analytic vanishing order
\lean{analyticOrderNatAt riemannZeta} at $\rho$. For real $\sigma,a,b$, set
\[
 Z(\sigma;a,b)=\{\rho:\zeta(\rho)=0,\ 0<\Rea\rho<1,\
                  \sigma\le\Rea\rho,\ a\le\Ima\rho\le b\}.
\]
Define the slab and density counts
\[
 n(\sigma;a,b)=\sum_{\rho\in Z(\sigma;a,b)}m(\rho),\qquad
 N(\sigma,T)=n(\sigma;-T,T).
\]
\end{definition}
\source{MathCollab.Density.ZetaCounting: IsNontrivialZetaZero, zetaMultiplicity, zetaZeroRegion, zetaZeroFinset, zetaSlabCount, zetaDensityCount}

\begin{lemma}[The finite-set and multiplicity conventions are genuine]
Every set $Z(\sigma;a,b)$ is finite. At a nontrivial zero, $m(\rho)$ is finite
and positive, and there is an analytic function $g$ near $\rho$, with
$g(\rho)\ne0$, such that locally
\[
 \zeta(z)=(z-\rho)^{m(\rho)}g(z).
\]
\end{lemma}
\begin{proof}
For $\rho\in Z(\sigma;a,b)$, the source bounds
$|\Ima\rho|\le |a|+|b|$ and
\[
 |\rho|\le |\Rea\rho|+|\Ima\rho|\le1+|a|+|b|.
\]
Thus the zero set is contained in the intersection of a closed ball and the
actual zeta-zero set. The imported mathlib theorem that this intersection is
finite gives the required finite set, without assuming that the open strip is
compact. Zeta is analytic away from $1$. If its analytic order at any $\rho\ne1$
were infinite, it would vanish on a neighborhood of $\rho$. The identity theorem
on the connected punctured plane $\C\setminus\{1\}$ would make it identically
zero there. This contradicts its nonvanishing at $2$. The finite-order local
factorization theorem for analytic functions then gives $g$. At a zero the
order is not zero, so its conversion to a natural number is strictly positive.
\end{proof}
\source{MathCollab.Density.ZetaCounting: zetaZeroRegion_finite, zeta_analyticOrder_ne_top, zetaMultiplicity_pos, zetaMultiplicity_local_factorization}

\begin{lemma}[Monotonicity, inclusions, and the right edge]
If $\sigma'\le\sigma$, $a'\le a$, and $b\le b'$, then
$n(\sigma;a,b)\le n(\sigma';a',b')$. Consequently if also $T\le T'$, then
$N(\sigma,T)\le N(\sigma',T')$. If $\sigma\ge1$, $N(\sigma,T)=0$.
The number of zero locations in a slab is at most its multiplicity count.
\end{lemma}
\begin{proof}
The inequalities on parameters give an inclusion of zero sets. Sum the
nonnegative natural multiplicities over that inclusion. If $\sigma\ge1$, a
putative zero has both $\Rea\rho\ge\sigma\ge1$ and $\Rea\rho<1$, a contradiction.
For the last assertion sum $1\le m(\rho)$ over the finite zero set.
\end{proof}
\source{MathCollab.Density.ZetaCounting: zetaSlabCount_mono, zetaDensityCount_mono, zetaZeroFinset_card_le_count; GuthMaynard.Improvement.GlobalDensityConditional: zetaDensityCount_zero_of_one_le}

\section{Source polynomials and separated observation sets}
\begin{definition}
For $N\in\R$ put
\[
 I_N=\{n\in\N:\lfloor N\rfloor_+<n\le\lfloor2N\rfloor_+\},\qquad
 D_{N,a}(t)=\sum_{n\in I_N}a_n\e^{\ii t\log(n/N)}.
\]
For $N\ge0$, membership in $I_N$ is equivalent to $N<n\le2N$.
A finite set $W\subset\R$ is one-separated if
$|t-u|\ge1$ whenever $t,u\in W$ and $t\ne u$.
\end{definition}
The coefficient hypothesis throughout is $|a_n|\le1$ for $n\in I_N$;
no hypothesis is placed on coefficients outside that interval.
\source{GuthMaynard.Improvement.Target: boundedPolynomial; SourceRatioAssembly: sourceRatioIndices_mem}

\begin{lemma}[Normalization and real-to-integer source transfer]
For $N>0$, the norm of $D_{N,a}$ equals the norm of
$\sum_{n\in I_N}a_n\e^{\ii t\log n}$. For $N\ge2$ and $M=\lfloor N\rfloor_+$,
\[
 M\ge2,\quad N/2\le M\le N,\quad
 \lfloor2N\rfloor_+\in\{2M,2M+1\}.
\]
For bounded source coefficients,
\[
 |D_{N,a}(t)|\le |D_{M,a}(t)|+1.
\]
In particular a value at least $V\ge2$ transfers to the integer-length
polynomial at threshold $V/2$.
\end{lemma}
\begin{proof}
Factor out $\e^{-\ii t\log N}$, which has modulus one. The floor inequalities
$M\le N<M+1$ give the bounds on $M$ and the two possible values of
$\lfloor2N\rfloor_+$. The integer sum at real length differs from the sum at
length $M$ by either no term or the term with index $2M+1$. That term has norm
at most one. The triangle inequality proves the comparison. Finally
$V-1\ge V/2$ if $V\ge2$.
\end{proof}
\source{GuthMaynard.Improvement.RealIntegerSource: source_floor_bounds, source_double_floor, boundedPolynomial_norm_integer_sum, boundedPolynomial_le_detecting_floor_add_one, detecting_floor_large_value}

\section{Endpoint claims, with the exact quantifier order}
Set
\[
 A(\sigma)=
 \begin{cases}
  3/(2-\sigma),&\sigma\le9/13,\\
  18/(7\sigma+3),&9/13<\sigma\le5/7,\\
  9/4,&5/7<\sigma\le3/4,\\
  3/(2\sigma),&3/4<\sigma,
 \end{cases}\qquad f(\sigma)=A(\sigma)(1-\sigma).
\]
Only the interval $1/2\le\sigma\le1$ is asserted in the endpoint density results.

\begin{theorem}[The fixed-line full curve]\label{thm:full}
For each $\sigma\in[1/2,1]$ and each $\eps>0$ there is
$C_{\sigma,\eps}>0$ such that for every real $T\ge2$,
\[
 N(\sigma,T)\le C_{\sigma,\eps}T^{f(\sigma)+\eps}.
\]
\end{theorem}
\begin{theorem}[The globally uniform coefficient]\label{thm:global}
For each $\eps>0$ there is $C_\eps>0$ such that simultaneously for
$\sigma\in[1/2,1]$ and $T\ge2$,
\[
 N(\sigma,T)\le C_\eps T^{(39/17)(1-\sigma)+\eps}.
\]
\end{theorem}
\begin{theorem}[Uniform exact curve away from the strict boundary]\label{thm:away}
For each $\delta>0$ and $\eps>0$ there is $C_{\delta,\eps}>0$ such that
for every $\sigma\in[1/2,1]$ satisfying
$\sigma\le3/4$ or $\sigma\ge3/4+\delta$, and every $T\ge2$,
\[
 N(\sigma,T)\le C_{\delta,\eps}T^{f(\sigma)+\eps}.
\]
\end{theorem}
\begin{theorem}[Middle-range large values]\label{thm:middle}
For each $\eps>0$ there is $C_\eps>0$ such that for all real
$N,T,V,t_0$, all coefficients $a:\N\to\C$, and all finite $W\subset\R$,
if
\[
 N\ge2,\quad N^{6/5}\le T\le N^{5/3},\quad V>0,
\]
$|a_n|\le1$ on $I_N$, $W\subset[t_0,t_0+T]$, $W$ is one-separated,
and $|D_{N,a}(t)|\ge V$ for $t\in W$, then
\[
 |W|\le C_\eps N^\eps
       \left(\frac{N^2}{V^2}+\frac{N^{18/7}T^{6/7}}{V^4}\right).
\]
\end{theorem}
\source{GuthMaynardStatement.full_curve, global_uniform, uniform_away_boundary, middle_large_values in lean/Solution.lean}
The proofs of Theorems~\ref{thm:full}--\ref{thm:away} appear after all branches
have been established. Theorem~\ref{thm:middle} is proved in Chapter~\ref{ch:middle}.

\chapter{Fourier-affine machinery and its uniform interfaces}\label{ch:affine}
\section{The exact affine average}
\begin{definition}
For $R,K\ge1$ and $C_0\ge0$, let
\[
 \mathcal B(R,K,C_0)=\{(b,q,d)\in\Z^3:
 R\le|b|\le2R,\ K\le q\le2K,\ |d|\le C_0R\}.
\]
The slope has both signs; the denominator is positive. Put
\[
 J_{R,K,C_0}(g)=\int_\R
 \left(\sum_{(b,q,d)\in\mathcal B(R,K,C_0)}
                  g\left(\frac{bu+d}{q}\right)\right)^2\,du.
\]
For a finite subset $S$ the corresponding energy $J_S(g)$ uses $S$ in the sum.
For a smooth compactly supported input, let
$\|g\|_\infty=\sup_x|g(x)|$.
A fixed Fourier profile $D(\beta,n)$ means the family of bounds
\begin{equation}\label{eq:supprofile}
 |\Hhat g(\xi)|\le D(\beta,n)\mathcal T^\beta
       (\mathcal T/|\xi|)^n\|g\|_\infty
 \quad(\beta>0,\ n\in\N,\ \xi\ne0).
\end{equation}
\end{definition}
\source{GuthMaynard.Improvement.StrengthenedAffine: sourceSup, FourierSupInput, affineAverage; AffineBox: InAffineBox; AffineFiniteCover: affineIndexIntegral}

\begin{analyticinput}[The strengthened unequal-scale affine bound]\label{input:affine}
Fix $0<a\le b$, $C_0\ge0$, and a nonnegative Fourier profile $D$.
For each $\eps>0$ there is $C>0$ such that for all
$\mathcal T\ge2$, $1\le K\le R\le\mathcal T^4$, and all nonnegative
$C^\infty$ functions $g$ supported in $[a,b]$ satisfying
\eqref{eq:supprofile},
\begin{equation}\label{eq:affine}
 J_{R,K,C_0}(g)\le C\mathcal T^\eps
 \left(R^2K^4\left(\int g\right)^2+RK^3\int g^2\right).
\end{equation}
The constant precedes the height, both real scales, and the function. There is
also a finite-subset and mixed-box version: slope and translation may have
separate shells bounded by $R$, and the denominator shell is bounded by $K$,
with the same right-hand expression up to the fixed support constants.
\end{analyticinput}
\source{GuthMaynard.Improvement.StrengthenedAffine.strengthened_unequal_affine; UnequalFullInput.unequal_affine_full_finite; MixedHeightTheorem.mixed_height_affine; MixedMaskedAffine.masked_mixed_affine}

\paragraph{How this input is proved in the development.}
The following explanation retains the actual two-stage reduction; it does not
replace this input by an unproved affine conjecture. For the source shells
$M_1,M_2,M_3$, the native affine transform integrates the square of
$\sum g((m_1u+m_3)/m_2)$; the first two shells are signed nonzero dyadic
integer shells and the third is a central integer shell. Its energy
$\mathcal J(g,M)$ is the finite maximum over $1\le M_j\le M$.
The Fourier-height family is mass-normalized:
\[
 \sup_\xi |\xi|^n|\Hhat g(\xi)|
 \le 2^dD(\beta,n)\mathcal T^\beta\mathcal T^{n+2}\int g
\]
at smoothing depth $d$. The support envelope at that depth is
$[1/4-2d/\mathcal T,\,9/4+2d/\mathcal T]$.
The factor $\mathcal T^{n+2}$ is not a height-independent constant.

The native Lemma-9.2 reduction supplies, after choosing its Fourier order,
a recurrence of the form
\begin{equation}\label{eq:affrec}
 \mathcal J(g,M)\le A\mathcal T^\alpha M^6\Bigl(\int g\Bigr)^2
 +B\mathcal T^\alpha M^2\|g\|_2
       \sqrt{\mathcal J(\widetilde g,M_2)}
 +E\Bigl(\int g\Bigr)^2,
 \qquad M_2\le M.
\end{equation}
Here $A,B,E$ are uniform functions of the selected Fourier order, divisor
constant, and current profile. The adaptive smoothing is the actual
\lean{gmAffineTildeSchwartz}; its parameter equals
$\mathcal T+M_2/(2M_3)$ after division by the first Poisson window.
It preserves nonnegativity, enlarges the support by at most $2/\mathcal T$,
and updates the Fourier coefficient by a factor two. The source proves
\[
 M_2^6\Bigl(\int\widetilde g\Bigr)^2+M_2^4\int\widetilde g^2
 \le16\left(M^6\Bigl(\int g\Bigr)^2+M^4\int g^2\right).
\]
Let the latter parenthesis be $W$. For a desired loss $\delta$, choose finitely
many descent steps $r$ with $(3/2)^r\delta\ge100$. At the terminal loss
$\ge100$, the crude source bound
$\mathcal J(g,M)\le36992M^6\int g^2$ and $M\le\mathcal T^4$ suffice.
At an induction step use \eqref{eq:affrec} with $\alpha=\delta/8$ and the
induction hypothesis with loss $3\delta/2$ on $\widetilde g$. It gives
\[
 \sqrt{\mathcal J(\widetilde g,M_2)}
 \le4\sqrt{C_{\rm next}}\mathcal T^{3\delta/4}\sqrt W.
\]
Since $M^2\|g\|_2\le\sqrt W$, the mixed term is at most
$4B\sqrt{C_{\rm next}}\mathcal T^{7\delta/8}W$.
The first and last terms are at most $(A+E)\mathcal T^\delta W$.
Thus $C=A+E+4B\sqrt{C_{\rm next}}+1$ closes the induction. A threshold
$\mathcal T\ge\max(2,8r)$ guarantees all support envelopes used by this
finite descent. This is why the complete profile must be fixed before the
varying function and height.
\source{RiemannZeta.GuthMaynard.exists_uniform_gmAffineJ_le_height_finite_descent and gmAffine_proposition9_1_height_family_native in GuthMaynard/AffineSource/LargeValuesAffineHeight.lean}

The unequal-scale Poisson reduction retains an outer factor $RK$, not a
factor $R^2$. With $m=\int g$ and $n=\|g\|_2$ it supplies
\[
 J\le E_0(R^2K^4m^2+RK n\sqrt Q),\qquad
 Q\le H_0(K^6m^2+K^4n^2),\quad H_0\ge1.
\]
Hence $\sqrt Q\le\sqrt{H_0}(K^3m+K^2n)$. The cross term
$RK^4mn$ is at most $R^2K^4m^2+RK^3n^2$, by expanding a nonnegative
square and using $K\le R$. Therefore
\[
 J\le2E_0\sqrt{H_0}(R^2K^4m^2+RK^3n^2).
\]
The source takes $E_0=A_0\mathcal T^{\eps/2}$ and
$H_0=(B_0+1)\mathcal T^\eps$, so the total loss is exactly
$\mathcal T^\eps$. Heights below $8$ are treated by the crude finite
bound and $R\le8^4$, not discarded.
\source{GuthMaynard.Improvement.UniformReduction.unequal_reduction_of_estimates; UnequalHeightTheorem.unequal_height_affine}

\paragraph{Exact Poisson normalizations.}
The identities used in that reduction are
\[
 \Hhat{(\psi(\cdot/R)e(\alpha\cdot))}(\xi)
       =R\Hhat\psi(R(\xi-\alpha)),\qquad
 \sum_{d\in\Z}\psi(d/R)e(\alpha d)
       =R\sum_{\ell\in\Z}\Hhat\psi(R(\ell-\alpha)).
\]
The second Poisson scale is $S/K$, so its argument is
$(S/K)(j-q'u'+qu)$ and its prefactor is $S/K$.
The affine Fourier factor for positive $b,q$ is
$e(\xi d/b)(q/b)\Hhat g(q\xi/b)$. Signed slopes are handled by
explicit reindexing. All interchanges require the source's integrability
and rapid-decay proofs. Passing from canonical support to general fixed
positive support uses $g_1(x)=g(\ell(x-1))$ and the injective index map
$(b,q,d)\mapsto(b,\ell q,d+\ell q)$; there is no Jacobian in $u$.
\source{GuthMaynard.Improvement.PoissonReduction; UnequalFullInput.supportIndexMap_term}

\section{Masks and smoothing without regularity of mask boundaries}
Let $b_\kappa$ be the source's \lean{ContDiffBump} at zero with inner
radius $1$ and outer radius $2$. Its normalized kernel is the exact function
$\kappa(x)=b_\kappa(x)/\int b_\kappa$; the denominator is strictly positive.
Thus $\kappa$ is nonnegative, compactly supported, smooth and has mass one.
Set $\kappa_h(x)=h^{-1}\kappa(x/h)$ and $S_hg=\kappa_h*g$, for $h>0$.
\source{GuthMaynard.Improvement.CompactKernel: kernelBump, compactKernel, compactKernel_mass}
\begin{lemma}[Mass and square norm under smoothing]
For nonnegative integrable $g$ with $g^2$ integrable,
\[
 \int S_hg=\int g,\qquad \int(S_hg)^2\le\int g^2.
\]
If $g$ is supported in $[a,b]$, $S_hg$ is smooth and supported in
$[a-2h,b+2h]$.
\end{lemma}
\begin{proof}
Tonelli and the mass-one identity give preservation of mass. Apply the
weighted Cauchy--Schwarz inequality at each $x$ to
$\int\kappa_h(y)g(x-y)\,dy$, then integrate and use translation invariance.
The support inclusion follows because $x-y$ cannot be in $[a,b]$ when
$x$ is outside its $2h$ enlargement and $|y|\le2h$. Derivatives may be taken
on the compactly supported smooth convolution factor, as justified by the
source's convolution regularity theorem.
\end{proof}
\source{GuthMaynard.Improvement.Smoothing: smooth_mass, smooth_l2_contraction; MaskedAffine.compact_smoothed_general_schwartz}

\begin{proposition}[The masked affine interface]\label{prop:maskedaff}
Fix $0<a\le b$ and $C_0\ge0$. For each $\eps>0$ one constant works for
all admissible $\mathcal T,R,K,h$, every measurable set $E$, and every
measurable real $r$ supported in $[a,b]$ with $r^2,r^4$ integrable:
\[
 J_{R,K,C_0}\bigl(S_h(\one_Er^2)\bigr)
 \le C\mathcal T^\eps
 \left(R^2K^4\Bigl(\int_Er^2\Bigr)^2+RK^3\int_Er^4\right),
\]
where $\mathcal T\ge2$, $1\le K\le R\le\mathcal T^4$, and
$0<h\le\min(1,a/4)$, $h^{-1}\le\mathcal T$.
\end{proposition}
\begin{proof}
The smoothed function is nonnegative, smooth, and supported in
$[a/2,b+2]$. The established normalized smoothing theorem supplies
\eqref{eq:supprofile} with constants independent of $E,r,h,\mathcal T$;
importantly the normalization is by the actual output supremum, not by an
inverse mask mass. Apply Input~\ref{input:affine} on the enlarged fixed
interval. The preceding lemma preserves $\int_Er^2$ and bounds the output
square integral by $\int_Er^4$. This proves the assertion including the
case of zero restricted mass. No regularity of $\partial E$ is used.
\end{proof}
\source{GuthMaynard.Improvement.MaskedAffine: compact_smoothed_full_input, compact_smoothed_unequal_affine, masked_unequal_affine}

\chapter{The centered cubic, exceptional Fourier mass, and spectral comparison}\label{ch:cubic}
\section{Literal kernels, measures, and operators}
Fix the actual comparison bump $w$: it is the smooth bump centered at $5/2$
with inner radius $3/2$ and outer radius $2$. Thus $0\le w\le1$,
$w=1$ on $[1,4]$, and its closed support is contained in $[1/2,5]$.
Put $a_0(x)=w(x)^2$ and $\chi(s)=w(s+2)$.
Then $0\le\chi\le1$, $\chi=1$ on $[-1,2]$, and
$\supp\chi\subset[-2,3]$.
\source{MathCollab.Density.ComparisonCutoff; GuthMaynard.Improvement.FixedCutoffs}

For a smooth amplitude $a$ supported in $[l,u]\subset(0,\infty)$ define
\begin{equation}\label{eq:kernel}
 K_{a,L}(s)=\sum_{\lceil lL\rceil\le m\le\lfloor uL\rfloor}
             a(m/L)\e^{\ii s\log(m/L)}
           -L\int_\R a(x)\e^{\ii s\log x}\,dx.
\end{equation}
For a measurable nonnegative integrable time density $\mu$, set
$d\nu_\mu(t)=\mu(t)\,dt$ and
\[
 \mathcal C_{a,L}(\mu)=\iiint
 K_{a,L}(t_1-t_2)K_{a,L}(t_2-t_3)K_{a,L}(t_3-t_1)
 \prod_{j=1}^3d\nu_\mu(t_j).
\]
This is the independent definition \lean{kernelCubic}, not a trace definition.
In $\mathcal H=L^2(\nu_\mu;\C)$, let
$p_x(t)=e(t\lambda(x))$. For a real nonnegative cutoff $a$, define
\[
 G=\sum_{\lceil lL\rceil\le m\le\lfloor uL\rfloor}
 a(m/L)\,p_{m/L}\otimes p_{m/L},\qquad
 C=L\int a(x)\,p_x\otimes p_x\,dx,\qquad H=G-C.
\]
Here $(v\otimes v)z=\langle v,z\rangle v$.
These are positive operators $G,C$, while $H$ need not be positive.
The source constructs the integral as a nuclear integral of rank-one operators;
$\int\|La(x)p_x\|\|p_x\|\,dx<\infty$ follows from $a\in L^1$ and
$\|p_x\|^2=\nu_\mu(\R)$.
\source{GuthMaynard.Improvement.OperatorKernel: operatorDiscrete, operatorContinuous; OperatorCubic; CubicKernel; CenteredPoisson}

\begin{lemma}[The genuine cubic trace identity]
For a countable Hilbert basis $b$ of $\mathcal H$,
\[
 \tr_b(H^3)=\mathcal C_{a,L}(\mu),\qquad
 \tr_b(T)=\sum_j\langle b_j,Tb_j\rangle
\]
for the nuclear operators used here.
\end{lemma}
\begin{proof}
Expand the nuclear presentations of the three factors. A rank-one cyclic
product has trace the cyclic product of three feature inner products.
Since $\langle p_x,p_y\rangle=\Hhat\mu(\lambda(x)-\lambda(y))$ with the
source's orientation, integration over frequencies gives the Fourier-side
cyclic product; Fubini gives the independently defined time-side kernel
cubic. The discrete-continuous difference is represented by the disjoint
union of finite counting frequencies and continuous frequencies, with signed
amplitude. Absolute integrability of the products is bounded by the product
of the nuclear norm integrals. This justifies the expansion and every
interchange. The source proves basis-independence through these integral
presentations, rather than presuming an infinite eigenbasis.
\end{proof}
\source{GuthMaynard.Improvement.OperatorCubic.operator_frequency_cubic; OperatorCentered.operator_centered_cubic; OperatorTrace}

\section{The masked cubic estimate and why the first mask is retained}
For the remainder of this section $\mu\ge0$ is measurable, belongs to
$L^1\cap L^2$, and is supported in an interval of length at most $\mathcal T$.
Write
\[
 F=\Hhat\mu,\quad M=\int\mu^2,\quad M_1=\int\mu,\quad
 E=\int|F|^4,\quad E_0=\int_{B^c}|F|^4,\quad A_2=\int_B|F|^2.
\]
The Fourier square and fourth powers are integrable: Plancherel gives
$\int|F|^2=M$, and $|F|\le M_1$ gives
$\int|F|^4\le M_1^2M$.

\begin{analyticinput}[Complete masked centered-cubic estimate]\label{input:maskedcubic}
For fixed smooth $a$ supported in a fixed positive interval, each $\eps>0$
has a constant $C>0$ such that if
$\mathcal T\ge2$, $\mathcal T^{4/5}\le L\le\mathcal T$, and $B$ is
measurable and symmetric under $\xi\mapsto-\xi$, then
\begin{equation}\label{eq:maskedcubic}
 |\mathcal C_{a,L}(\mu)|\le C\mathcal T^\eps
 \left(\mathcal T^2M^{3/2}+\mathcal TL\sqrt M\sqrt{E_0}
    +\sqrt{A_2}(\mathcal T^2M+\mathcal TL\sqrt E)\right).
\end{equation}
The constant precedes all densities, masks, and scales.
\end{analyticinput}
\source{GuthMaynard.Improvement.CenteredRestrictedCubic.centered_masked_cubic_estimate}

\paragraph{Precise analytic construction behind this input.}
Poisson summation turns \eqref{eq:kernel} into the sum of the nonzero modes
$L\Hhat{(a(x)\e^{\ii s\log x})}(Lm)$. The source retains
$0<|m|\le U$, $U=\mathcal T^{1+\eta}/L$, choosing $0<\eta<4/5$ before
all varying data. It bounds the cubic replacement error by
$C\mathcal T^{-100}M^{3/2}$. Expanding the retained cube gives the exact
factor $L^3$ and, for $m=(m_1,m_2,m_3)$, the spatial integral with amplitude
$a(x_1)a(x_2)a(x_3)$, phase $e(-L(m_1x_1+m_2x_2+m_3x_3))$, and Fourier product
\[
 F(\lambda(x_3)-\lambda(x_1))
 F(\lambda(x_1)-\lambda(x_2))
 F(\lambda(x_2)-\lambda(x_3)).
\]
Set $x_1=rv$, $x_2=ry$, $x_3=r$. The Jacobian is exactly $r^2$.
The Fourier product becomes
$F(-\lambda(v))F(\lambda(v)-\lambda(y))F(\lambda(y))$.
Radial integration, before taking a norm, yields rapid decay in
$L(m_1v+m_2y+m_3)$. The near region is the thin strip
$|L(m_1v+m_2y+m_3)|\le\mathcal T^\eta$.
It is rewritten as a window of width
$\mathcal T^\eta/(L|m_2|)$ about
$-(m_1v+m_3)/m_2$. Negative denominators are normalized by simultaneous
negation of modes. For the inverse-ratio energy, the exact identity
\[
 \frac{-(m_1v+m_3)/m_2}{v}
 =-\frac{m_3/v+m_1}{m_2}
\]
swaps slope and translation. The Jacobian for $v\mapsto1/v$ is bounded on
the fixed positive ratio interval. Both energies are bounded by the mixed
masked affine theorem of Chapter~\ref{ch:affine}. The dyadic mode sum,
width losses, mode count, and radial discard are charged in the source
\lean{cubic_scale_closure} chain. In Fourier coordinates the relevant
measure change is $dx=2\pi x\,d\lambda$; no unrecorded normalization change
is made.

For a symmetric exceptional set $B$, the retained cubic is partitioned exactly
into four designated-mask terms
\[
 (B^c,B^c,B^c),\qquad(B,\R,\R),\qquad
 (B,\R,B^c),\qquad(B,B^c,B^c).
\]
Cyclic reindexing of both spatial coordinates and the full mode cube moves
each exceptional mask to the first Fourier factor. The all-low term costs
$\mathcal T^\eps(\mathcal T^2M^{3/2}+\mathcal TL\sqrt M\sqrt{E_0})$;
each of the other three costs
$\mathcal T^\eps\sqrt{A_2}(\mathcal T^2M+\mathcal TL\sqrt E)$.
The homogeneous negligible errors are absorbed by the first term.
For bounded heights the source separately bounds the cubic by a constant
multiple of $\mathcal T^2M^{3/2}$. Thus eventual Poisson retention is not
used to skip the bounded-height range.
\source{GuthMaynard.Improvement.CubicRetainedPartition.retained_cubic_designated_partition; RadialCoordinates; CubicMixedBox; CubicFourierMoments; CubicMaskedEstimate; CubicRetainedMasked; CenteredCubicApprox; CenteredCubicBoundedHeight}

\paragraph{Explicit affine-to-cubic exponent closure.}
For fixed ratio support and a fixed factor $q_0\ge1$, the smoothing width
in a denominator shell $K$ is $h=q_0\mathcal T^\eta/(LK)$.
For $1\le K\le R\le U=\mathcal T^{1+\eta}/L$ one has
$h^{-1}\le\mathcal T$ and
$h\le q_0\mathcal T^{\eta-4/5}$. Since $\eta<4/5$, a single eventual
height threshold puts every such width in the admissible compact-smoothing
range of Proposition~\ref{prop:maskedaff}.
After the affine estimate and Cauchy--Schwarz, a box contribution at restricted
first-factor square mass $A$ and other-factor masses $M_D,E_D$ has scalar factor
\[
 \sqrt A\frac{L^2\mathcal T^\eta}{K}
       \sqrt{R^2K^4M_D^2+RK^3E_D}.
\]
The square root is at most $K(RKM_D+\sqrt{RK}\sqrt{E_D})$.
Both $RK\le U^2$ and $\sqrt{RK}\le U$; hence the factor is at most
\[
 \sqrt A\left(\mathcal T^{3\eta}\mathcal T^2M_D
                +\mathcal T^{2\eta}\mathcal TL\sqrt{E_D}\right),
\]
because $L^2\mathcal T^\eta U^2=\mathcal T^{2+3\eta}$ and
$L^2\mathcal T^\eta U=\mathcal T^{1+2\eta}L$.
A Fourier-affine loss $\mathcal T^\alpha$ contributes
$\mathcal T^{\alpha/2}$ after the square root. The finite dyadic-class
count contributes at most $D_\delta\mathcal T^\delta$, since $U\le\mathcal T$.
Choose $\eta=\min(\eps/12,1/10)$, $\alpha=\eps/4$, $\delta=\eps/4$;
then $\delta+\alpha/2+3\eta\le\eps$.
Thus both terms fit the displayed masked-cubic budget while the restricted
factor $\sqrt A$ is retained unchanged. Passing back from ratio coordinates
charges only fixed factors bounded using $dx=2\pi x\,d\lambda$.
\source{GuthMaynard.Improvement.CubicScaleClosure: cubic_smoothing_width_bounds, cubic_box_polynomial_scale, cubic_box_epsilon_absorption; CubicScaledStrip; CubicFourierMoments}

\begin{proposition}[Restricted cubic from the masked estimate]\label{prop:restricted}
For the same $a,\mu,\mathcal T,L$ and every $\Lambda\ge0$,
\begin{equation}\label{eq:restricted}
 |\mathcal C_{a,L}(\mu)|\le C\mathcal T^\eps
  \left(\mathcal T^2M^{3/2}+\mathcal TL\Lambda M
       +\mathcal TL M_1\int_{\{|F|>\Lambda\}}|F|^2\right).
\end{equation}
\end{proposition}
\begin{proof}
Take $B=\{|F|>\Lambda\}$. The Fourier transform is continuous, so $B$ is
measurable. Because $\mu$ is real, $F(-\xi)=\overline{F(\xi)}$, so $B$ is
symmetric. Put $A_2=\int_B|F|^2$. Plancherel and the pointwise bounds give
\[
 A_2\le M,\qquad E_0\le\Lambda^2M,\qquad
 E\le\Lambda^2M+M_1^2A_2.
\]
It follows that $\mathcal TL\sqrt M\sqrt{E_0}\le\mathcal TL\Lambda M$
and $\mathcal T^2M\sqrt{A_2}\le\mathcal T^2M^{3/2}$.
Also $\sqrt E\le\Lambda\sqrt M+M_1\sqrt{A_2}$; hence
\[
 \mathcal TL\sqrt{A_2}\sqrt E
 \le\mathcal TL\Lambda M+\mathcal TL M_1A_2.
\]
Substituting these three inequalities in \eqref{eq:maskedcubic} increases its
constant by at most a factor three. This is precisely the source's restricted
closure, including zero mass and zero threshold.
\end{proof}
\source{GuthMaynard.Improvement.CubicOrdinaryClosure.ordinary_fourier_threshold_budgets, ordinary_restricted_cubic_of_masked; CenteredRestrictedCubic.centered_restricted_cubic}

\section{The amplifier and the spectral inequality}
The amplifier column set is
\[
 Q_{N,L}=\{q\in\N:\lceil L/N\rceil_+\le q\le\lfloor2L/N\rfloor_+\}.
\]
For $N>0$, $N\le L$, its cardinality $k$ satisfies
$L/(2N)\le k\le2L/N$. Its divisor matrix is
\[
 B_{m,q}=\sum_{n\in I_N}\one_{m=qn}a_n.
\]
Every contributing product $qn$ has $1\le qn/L\le4$, the plateau of $w$.
Its synthesized column is exactly
$\e^{\ii t\log(qN/L)}D_{N,a}(t)$.
The source bounds the amplifier norm by divisor multiplicity, using the
established subpower divisor bound $d(n)\ll_\eta n^\eta$.
It then constructs a finite orthonormal family in its synthesized range,
not a postulated infinite eigenfamily.
\source{GuthMaynard.Improvement.AmplifierArithmetic; AmplifierNorm; AmplifierSynthesis; AmplifierQuantitative}

\begin{analyticinput}[The actual finite amplifier witness]\label{input:amp}
For each $\eta>0$ there is $P>0$ such that, uniformly in the bounded
coefficients, scales $N>0$, $N\le L$, and finite nonnegative time measures,
there is an orthonormal family $v_1,\ldots,v_r$, $r\le k$, satisfying
\[
 k\frac{\int\mu|D_{N,a}|^2}{P(1+L)^\eta(N+1)}
 \le\sum_{j=1}^r\langle v_j,Gv_j\rangle.
\]
Here $G$ uses amplitude $w^2$ on the indicated integer rows.
The source also proves the Schur bound $\|C\|\le D_0LB$ whenever
$0\le\mu\le B$, with $D_0$ depending only on the fixed cutoff.
\end{analyticinput}
\source{GuthMaynard.Improvement.AmplifierQuantitative.amplifier_time_density_witness; OperatorSchur.operator_fixed_cutoff_schur}
\begin{proof}[Proof of the finite amplifier witness]
Let $\mathcal A$ be the literal divisor matrix between the finite column
and row Euclidean spaces. For a row $m$, the product fiber is
$F_m=\{(q,n)\in Q_{N,L}\times I_N:qn=m\}$.
Cauchy--Schwarz and $|a_n|\le1$ give
$|\mathcal A z(m)|^2\le|F_m|\sum_{(q,n)\in F_m}|z_q|^2$.
A nonempty fiber injects into the divisors of $|m|$ and has $0<m\le4L$.
Thus $|F_m|\le C_\eta(1+L)^\eta$. Summing over rows counts each $(q,n)$
at most once and gives
$\|\mathcal A\|^2\le P(1+L)^\eta(N+1)$.
Let $\mathcal B$ synthesize a row coordinate into $w(m/L)p_{m/L}$.
Then $\mathcal B\mathcal B^*=G$ exactly. Every synthesized column
$\mathcal B\mathcal A e_q$ has square norm $\int\mu|D|^2$ by the
plateau identity and its unit phase.
Choose an orthonormal basis $v_1,\ldots,v_r$ of the finite range of
$\mathcal B\mathcal A$; $r\le k$. Parseval in that range and adjoint
Parseval in the finite column space give
\[
 \sum_q\|\mathcal B\mathcal A e_q\|^2
 =\sum_j\|\mathcal A^*\mathcal B^*v_j\|^2
 \le\|\mathcal A\|^2\sum_j\|\mathcal B^*v_j\|^2
 =\|\mathcal A\|^2\sum_j\langle v_j,Gv_j\rangle.
\]
Insert the matrix bound and divide by its positive uniform upper bound.
This proves the displayed witness, including rank zero. The Schur estimate
is the separate proved continuous-kernel norm input stated above.
\end{proof}
\source{GuthMaynard.Improvement.AmplifierNorm.amplifier_matrix_energy; AmplifierWitness.amplifier_range_witness; AmplifierSynthesis.amplifier_synthesis_factorization}

\begin{lemma}[Finite spectral sign control]
If $C\ge0$ and $G=C+H\ge0$ in finite dimension, with eigenvalues $h_j$
of $H$, then
\[
 \sum_j(-h_j)_+^3\le\tr(CH^2),\qquad
 \sum_j(h_j)_+^3\le|\tr(H^3)|+\tr(CH^2).
\]
\end{lemma}
\begin{proof}
In an orthonormal eigenbasis of $H$ put $c_j=\langle b_j,Cb_j\rangle$.
Positivity gives $c_j\ge0$ and $c_j+h_j\ge0$. If $h_j\le0$,
$(-h_j)^3\le c_jh_j^2$; otherwise the left side is zero. Sum these
inequalities. Finally $(h_j)_+^3=h_j^3+(-h_j)_+^3$ and
$\sum h_j^3\le|\sum h_j^3|$. Notice that $C$ is not assumed to commute
with $H$ or to be diagonal.
\end{proof}
\source{GuthMaynard.Improvement.FiniteSpectral.finite_negative_trace, finite_positive_trace; OperatorSpectralComparison.operator_eigenbasis_diagonal_nonneg}

\begin{analyticinput}[Nuclear spectral comparison]\label{input:spectral}
For positive $G,C$ with the source's nuclear presentations, a finite
orthonormal witness satisfying
$ka\le\sum_j\langle v_j,Gv_j\rangle$, $r\le k$, and $\|C\|\le B$ obeys
\[
 k(a-B)_+^3\le|\tr_b((G-C)^3)|+\Rea\tr_b(C(G-C)^2).
\]
The countable basis $b$ exists for the actual $L^2$ time space.
\end{analyticinput}
\paragraph{Justification of the imported operator step.}
In finite dimension the source constructs an eigenbasis of $G-C$. Bessel's
inequality bounds the projections of the finite witness. The positive-part
cubic inequality follows from the preceding sign-control lemma together
with the finite witness's quadratic lower bound. For the actual space take
finite-dimensional spans of the witness and finitely many basis vectors,
compress both operators, apply the finite inequality, and pass to the
directed limit. Nuclear presentations prove convergence of both cubic and
mixed diagonal traces. Thus no infinite spectral-basis assumption or
unproved compact trace identity is required.
\source{GuthMaynard.Improvement.OperatorSpectralComparison.operator_finite_spectral_comparison; OperatorNuclearSpectral.operator_nuclear_spectral_comparison; OperatorFiniteTrace; OperatorTraceConvergence}

\begin{analyticinput}[The mixed source correction]\label{input:mixed}
Fix $A\ge1$ and $p>0$. For every $\eps>0$ one constant works whenever
$N\ge2$, $N\le L$, $0\le\mathcal T\le AN^p$, $\mu\in L^1\cap L^2$
is measurable and nonnegative, supported in an interval of length
$\mathcal T$, and $\mu\le|D_{N,a}|^2$:
\[
 |\tr_b(CH^2)|\le R_0N^\eps L^2N\mathcal T\int\mu^2.
\]
\end{analyticinput}
\source{GuthMaynard.Improvement.OperatorSourceMixed.operator_source_mixed}
The source discharges this with the following literal row and cycle bounds.
For any fixed reflection cutoff $v$, put
$c_v(s)=\int v(x)\e^{\ii s\log x}dx$ and $K=K_{v,L}$.
The established source-row theorem states, uniformly for $u$ on the time support,
\[
 \int\mu(t)|K(u-t)|^2dt\le R_vN^\eps LN\mathcal T.
\]
Its proof uses the actual reflected Mellin kernel: near rows use bounded
coefficient/product energy, signed shells use Fourier-reflected band estimates,
and finite dyadic summation includes negative reflection. It applies to the
genuine centered row, not to an unrelated absolute kernel.
The continuous correction has kernel $Lc_v$. The source proves
$\psi=|c_v|\in L^1$, bounded, by compact-smooth Mellin decay.
By the nuclear mixed trace identity its cycle is
\[
 \iiint Lc_v(u-v)K(v-t)K(t-u)\mu(u)\mu(v)\mu(t)\,du\,dv\,dt.
\]
Use $|K(-s)|=|K(s)|$ and
$|K(v-t)K(t-u)|\le(|K(v-t)|^2+|K(u-t)|^2)/2$.
The row bound controls the inner $t$ integral by $R_vN^\eps LN\mathcal T$.
The remaining double integral of $\psi(u-v)\mu(u)\mu(v)$ is at most
$\|\psi\|_1\int\mu^2$, by the source's $L^1$--$L^2$ convolution inequality.
This proves the mixed estimate with constant
$(1+\|\psi\|_1)R_v$. In the raw spectral application $v=a_0=w^2$,
which is itself an admissible reflection cutoff; hence the displayed
$G,C,H$ are exactly the ones in the row estimate.
\source{GuthMaynard.Improvement.ReflectionSourceRow.centered_source_row; MixedKernelRow.kernelMixed_time_row_bound; MixedCutoffEnvelope; OperatorMixedBound; SourceSpectralRaw}

\begin{proposition}[Raw spectral mass estimate]\label{prop:spectralraw}
Under the hypotheses of Input~\ref{input:mixed} and also
$\mu\le\min(|D_{N,a}|^2,B)$, $B\ge0$, let $M=\int\mu^2$.
For any $\eta,\eps>0$ the source supplies constants $P,D_0,R_0>0$ such that
\begin{equation}\label{eq:spectralraw}
 M^3\le4\bigl(P(1+L)^\eta(N+1)\bigr)^3
 \left((D_0LB)^3+\frac{2N}{L}
   \left(|\mathcal C_{a_0,L}(\mu)|+R_0N^\eps L^2N\mathcal T M\right)\right).
\end{equation}
\end{proposition}
\begin{proof}
Pointwise $\mu^2\le\mu|D|^2$, with both integrals justified by boundedness
of the finite polynomial. The amplifier witness therefore has level at
least $M/P_1$, where $P_1=P(1+L)^\eta(N+1)>0$.
Use Input~\ref{input:spectral}, $k\ge L/(2N)$, the Schur bound
$\|C\|\le D_0LB$, the cubic trace identity, and Input~\ref{input:mixed}.
Writing $z=(M/P_1-D_0LB)_+$ gives
\[
 \frac L{2N}z^3\le|\mathcal C|+R_0N^\eps L^2N\mathcal T M.
\]
Since $M/P_1\le D_0LB+z$ and
$(x+y)^3\le4(x^3+y^3)$ for nonnegative $x,y$, multiplication by $P_1^3$
gives exactly \eqref{eq:spectralraw}. This proof also works when $M=0$.
\end{proof}
\source{GuthMaynard.Improvement.SpectralMassAssembly.actual_spectral_mass_bound, spectral_mass_scalar; SourceSpectralRaw.source_spectral_raw}

\begin{proposition}[Uniformly absorbed spectral scale]\label{prop:spectralscale}
Fix $A\ge1$, $p>0$. For every $\delta>0$, one constant $S>0$ works for
$N\ge2$, $N\le L\le AN^p$, $0\le\mathcal T\le AN^p$, and the preceding
source densities:
\begin{equation}\label{eq:spectralscale}
 M^3\le SN^\delta\left((NLB)^3+\frac{N^4}{L}|\mathcal C_{a_0,L}(\mu)|
                         +N^5L\mathcal T M\right).
\end{equation}
\end{proposition}
\begin{proof}
Take $\eta=\delta/(6p)$ and the mixed-source loss $\eps=\delta/2$.
Then $3p\eta+\eps=\delta$. From
$1+L\le(1+A)N^p$ and $N+1\le2N$,
\[
 4\bigl(P(1+L)^\eta(N+1)\bigr)^3N^\eps
 \le32P^3(1+A)^{3\eta}N^{\delta+3}.
\]
Insert an additional $N^\eps\ge1$ in the nonnegative terms of
\eqref{eq:spectralraw} where necessary. Expanding $N^3$ into the three
terms yields $(NLB)^3$, $2N^4|\mathcal C|/L$, and
$2R_0N^5L\mathcal TM$. A uniform constant
$32P^3(1+A)^{3\eta}(D_0^3+2+2R_0)$ dominates their coefficients.
\end{proof}
\source{GuthMaynard.Improvement.SpectralScaleBudget: spectral_prefactor_budget, spectral_raw_scale_budget, source_spectral_scale}

\chapter{The literal truncation and the optimized fourth moment}\label{ch:moment}
\section{Truncation is not clipping}
For $N\ge2$, $T>0$, a bounded source $D=D_{N,a}$, and real $U$, put
\[
 \beta_T(t)=\chi(t/T),\qquad
 \mu_U(t)=\one_{\{|D(t)|\le U\}}\beta_T(t)|D(t)|^2,\qquad
 r_U(t)=\one_{\{|D(t)|>U\}}\beta_T(t)|D(t)|^2.
\]
The measure density is \emph{discarded} above $U$; it is not
$\min(|D|^2,U^2)$. Let
\[
 M=\int\mu_U^2,\qquad H_{\rm rem}=\int r_U,\qquad M_1=\int\mu_U.
\]
All these integrals are finite and nonnegative. Indeed $D$ is a finite
continuous sum with $|D|\le N+1\le2N$ and $\beta_T$ is bounded and
compactly supported. The source proves measurability and integrability of
$\mu_U$, $\mu_U^2$, and $r_U$, support in $[-2T,3T]$, and
\[
 0\le\mu_U(t)\le\min(|D(t)|^2,U^2),\qquad
 \mu_U+r_U=\beta_T|D|^2.
\]
For negative $U$, $\mu_U=0$; the domination statement remains valid.
\source{GuthMaynard.Improvement.SourceTailConditional.source_density_regular; TruncatedSourceAttachment.fixed_truncated_source_regular}

\section{The arithmetic sampling interface and the exceptional Fourier tail}
\begin{analyticinput}[Actual source sampling]\label{input:sampling}
For the fixed $\chi$, one $D_0>0$ works for all $N,T\ge2$ and bounded
source coefficients. For every complex test $h$ with $h\in L^1\cap L^2$,
$\Hhat h\in L^1$, and $\supp h\subset\supp\beta_T$,
\[
 \left|\int |D(t)|^2h(t)\,dt\right|
 \le D_0N^2\int|\Hhat h|+
      D_0N\sqrt{T\log(2N)}\left(\int|h|^2\right)^{1/2}.
\]
\end{analyticinput}
\source{GuthMaynard.Improvement.SourceDirichletTail.boundedPolynomial_source_sampling; SourceTailConditional.SourceSamplingBound}

\paragraph{Arithmetic content of this input.}
Expansion of $|D|^2$ reduces pairing to Fourier samples at
$\log(n/m)/(2\pi)$, $n,m\in I_N$. Reduce $(n,m)$ by its greatest common
divisor to $(a,b)$; its actual multiplicity is the number of pairs on that
primitive rational ray. It is at most $\lfloor2N\rfloor/\max(a,b)$.
Distinct primitive pairs of height at most $Y$ have logarithms separated
by at least $Y^{-2}$: their integer cross-products differ by at least one,
and $\log p-\log q\ge1/p$ for $p\ge q+1$.
At height $q$ there are at most $2q$ positive pairs. Thus squared
multiplicities through height $Y$ have sum at most
$2\lfloor2N\rfloor^2(1+\log Y)$.
The source decomposes exact low height plus half-open dyadic height shells,
uses disjoint Fourier-sampling neighborhoods and the proved derivative
variation estimate, and obtains the displayed sampling bound. In particular
it does not discard repeated ratios, the diagonal ratio $(1,1)$, or the
nonintegral endpoint term.
More precisely, set $Y=\min(2N,\max(2,\sqrt T))$ and choose the least $J$
with $2N\le2^JY$. Then $Y\ge2$, $Y^2\le2T$,
$\sum_{j<J}2^jY\le4N$ and $\sum_{j<J}(2^jY)^{-1}\le2/Y$.
If $Y=2N$ then $J=0$; otherwise $Y\ge\sqrt T$.
For a test supported in $[-cT,cT]$, with
$M_h=\int|h|^2$ and $I_h=\int_{-2}^2|\widehat h|$, the exact low-height
sampling cost is at most
\[
 \sqrt{16N^2\log(2N)}\sqrt{(2\pi Y^2+4\pi cT)M_h}.
\]
The established shell-variation bound at height $Q=2^jY$ is
$16\pi NQI_h+8\pi NcTQ^{-1}\sqrt{M_h}$.
The geometric sums therefore give high cost at most
$64\pi N^2I_h+16\pi c(NT/Y)\sqrt{M_h}$, with the second term absent
when $J=0$. Since $NT/Y\le N\sqrt T$ when it is present, and
$\log(2N)\ge1$, the low and high terms fit
\[
 (64\pi+8\sqrt{\pi(1+c)}+16\pi c)
      \left(N^2I_h+N\sqrt{T\log(2N)}\sqrt{M_h}\right).
\]
For the fixed time support take $c=3$ and enlarge the band integral to the
full Fourier $L^1$ integral. The finite expansion of the polynomial pairing
then gives Input~\ref{input:sampling}, with bounded coefficients applied
only inside $I_N$.
\source{GuthMaynard.Improvement.ReducedRatioGeometry; SourceRatioAssembly; SamplingVariation; SamplingFourier; SourceRatioHeightBudget; SourceRatioBound}

\begin{proposition}[Source Fourier mass above the actual removal threshold]\label{prop:tail}
There is a fixed $D_1>0$ such that if
$\Lambda\ge D_1(N^2+H_{\rm rem})$, then
\begin{equation}\label{eq:tail}
 \int_{\{|\Hhat\mu_U|>\Lambda\}}|\Hhat\mu_U|^2
       \le D_1N^2T\log(2N).
\end{equation}
\end{proposition}
\begin{proof}
Set $E=\{|\Hhat\mu_U|>\Lambda\}$, $A_2=\int_E|\Hhat\mu_U|^2$, and
$g=\mathcal F^{-1}(\one_E\Hhat\mu_U)$. Since $\Lambda>0$,
\[
 \int_E|\Hhat\mu_U|\le A_2/\Lambda,\quad
 |g(t)|\le A_2/\Lambda,\quad\int|g|^2=A_2.
\]
The source proves the inverse-Fourier identity
$\int\mu_Ug=A_2$; reality makes this bilinear identity have the stated
orientation, by the symmetric Fourier mask. For
$h=\beta_Tg$, its Fourier $L^1$ norm is bounded by
$J A_2/\Lambda$, $J=\int|\Hhat\chi|$, and its square norm by $A_2$.
All test regularity follows from the source's smooth cutoff convolution
lemmas. The decomposition $\beta_T|D|^2=\mu_U+r_U$ and the sampling
input imply
\[
 A_2\le (D_0N^2J+H_{\rm rem})A_2/\Lambda
           +D_0N\sqrt{T\log(2N)}\sqrt{A_2}.
\]
Choose $D_1$ large enough that
$2(D_0N^2J+H_{\rm rem})\le\Lambda$ and $4D_0^2\le D_1$.
Absorb the first term, obtaining
$A_2/2\le D_0N\sqrt{T\log(2N)}\sqrt{A_2}$.
If $A_2=0$ the result is immediate; otherwise divide by $\sqrt{A_2}$ and
square. This proves \eqref{eq:tail} without assuming that the removed mass
is small.
\end{proof}
\source{GuthMaynard.Improvement.SourceTailConditional.source_tail_of_sampling, source_tail_scalar_absorption; SourceDirichletTail.boundedPolynomial_source_tail}

\begin{analyticinput}[Actual second moments and classical level volume]\label{input:classical}
For $N\ge2$, $T\ge0$, bounded source coefficients, and one-separated
$W\subset[t_0,t_0+T]$,
\[
 \sum_{t\in W}|D(t)|^2\le102N(T+N\log(2N)).
\]
For $N\le T$ the source gives the simpler bound
$204TN\log(2N)$.
For each $\eps>0$, whenever $N\ge2$, $N\le T\le N^2$, $V>0$,
\begin{equation}\label{eq:classical}
 |W|\le C_\eps N^\eps\left(N^2/V^2+TN^4/V^6\right)
\end{equation}
for every one-separated source large-value set. The corresponding Lebesgue
level volume on an interval of length $qT$ is at most
$qC_\eps N^\eps(N^2/V^2+TN^4/V^6)$ for integer $q\ge0$.
Finally, for our truncated source with $N\le T$,
\[
 M_1\le612TN\log(2N).
\]
\end{analyticinput}
\source{GuthMaynard.Improvement.SeparatedSecondMoment; ClassicalLargeValues.classical_large_values; ClassicalLevelVolume.boundedPolynomial_classical_level_volume; TruncatedSourceMeanValue.fixed_truncated_source_mass}

\paragraph{No missing threshold case in the classical estimate.}
The high-threshold input behind \eqref{eq:classical} is the native integer
large-values theorem on $T^{1/3}\le M\le T$, with threshold
$V\ge M^{3/4+\kappa}$, giving
$T^\eta(M^2/V^2+M^3\sqrt T/V^4)$.
The real-to-integer transfer of Chapter~1 applies above
$16N^{3/4+\eps/8}$; nonempty $W$ implies $N\ge8$, which gives the integer
scale condition. Its second term fits \eqref{eq:classical}, since
\[
 N^3\sqrt T/V^4\le N^2/V^2+TN^4/V^6
\]
by expanding $(NV^2-N^2\sqrt T)^2\ge0$.
Below the threshold the separated second moment is used, with
$V^4\le16^4N^{3+\eps/2}$ and
$\log(2N)\le(\log2+2/\eps)N^{\eps/2}$.
Empty sets are handled directly. Thus the classical input is not applied
outside its large-value height range.
\source{GuthMaynard.Improvement.ClassicalLargeValues: classical_cross_budget, classical_low_scalar, classical_floor_scale, classical_large_values}

\begin{proposition}[Removed mass bound]\label{prop:removed}
For each $\eps>0$ one constant works for $N\ge2$, $N\le T\le N^2$,
$U\ge1$, and every bounded coefficient family:
\begin{equation}\label{eq:removed}
 H_{\rm rem}\le C_\eps N^\eps(N^2+TN^4/U^4).
\end{equation}
\end{proposition}
\begin{proof}
Choose a dyadic cutoff $J$ with $2N\le2^J$ and
$J+1\le K_\eps N^{\eps/2}$. Such a cutoff is obtained from the least
dyadic cover and the logarithmic power bound. Since $|D|\le2N\le2^JU$,
split the removed part into levels $2^jU<|D|\le2^{j+1}U$, $0\le j<J$.
On each level $|D|^2\le4(2^jU)^2$. The cutoff is at most one and supported
in $[-2T,3T]$, an interval of length $5T$, so Input~\ref{input:classical}
with loss $\eps/2$ gives
\[
 H_{\rm rem}\le5C N^{\eps/2}\sum_{j<J}
 4(2^jU)^2\left(N^2/(2^jU)^2+TN^4/(2^jU)^6\right).
\]
The first part is $4JN^2$. The second is
$4TN^4U^{-4}\sum_{j<J}16^{-j}\le8TN^4U^{-4}$.
Use the bound on $J+1$ and $N^{\eps/2}\ge1$; the product of the two
half-losses is $N^\eps$. The source chooses $40CK$ as a common constant.
\end{proof}
\source{GuthMaynard.Improvement.RemovedMassBound: removed_dyadic_classical_sum, removed_dyadic_cover, fixed_source_removed_mass}

\section{Fixed-source raw cubic and spectral assembly}
Use Proposition~\ref{prop:restricted} on support length $5T$, with
$\Lambda=D_1(N^2+H_{\rm rem})$. By \eqref{eq:tail} and
$M_1\le612TN\log(2N)$,
\begin{equation}\label{eq:fixedcubic}
 |\mathcal C_{a_0,L}(\mu_U)|\le C(5T)^\eps
 \left(25T^2M^{3/2}+5TLD_1(N^2+H_{\rm rem})M
         +3060D_1N^3LT^3\log^2(2N)\right).
\end{equation}
This requires $(5T)^{4/5}\le L\le5T$, not merely $T^{4/5}\le L$.
The coefficient $3060$ is the product $5\cdot612$, and the support
translation is $[-2T,3T]$.
\source{GuthMaynard.Improvement.TruncatedSourceCubic.fixed_truncated_source_cubic}

If also $N\le T\le N^2$, $N\le L$, then
$L,5T\le5N^2$. Apply Proposition~\ref{prop:spectralscale} with $A=5$,
$p=2$, $B=U^2$. Substituting \eqref{eq:fixedcubic} gives
\begin{align}\label{eq:fixedraw}
 M^3\le SN^\delta\bigg[&(NLU^2)^3+5N^5LTM\\
 &+\frac{N^4}{L}C(5T)^\eps
 \bigg(25T^2M^{3/2}+5TLD_1(N^2+H_{\rm rem})M\\
 &\hspace{50mm}+3060D_1N^3LT^3\log^2(2N)\bigg)\bigg].\nonumber
\end{align}
The constants $S,C,D_1$ precede all sources and scales.
\source{GuthMaynard.Improvement.TruncatedSourceMomentRaw.fixed_truncated_source_moment_raw}

\section{The three optimized quantities}
\begin{definition}
For $N,T>0$ define
\[
 L_*=N^{1/7}T^{5/7},\qquad U_*=N^{5/7}T^{1/14},\qquad
 P_*=N^{18/7}T^{6/7}.
\]
\end{definition}
\begin{lemma}[Admissibility and exact scale identities]\label{lem:scales}
If $N\ge2$ and $N^{6/5}\le T\le N^{5/3}$, then
\[
 N\le T\le N^2,\quad N\le L_*,\quad T^{4/5}\le L_*\le T,
 \quad1\le U_*\le N,
\]
and $5L_*$ is admissible for \eqref{eq:fixedraw}. Moreover
\begin{align*}
 NL_*U_*^2&=P_*,&N^4T^2/L_*&=P_*^{3/2}=P_*\sqrt{P_*},\\
 N^5L_*T&=P_*^2,&TN^4/U_*^2&=P_*,\\
 TN^4/U_*^4&=NL_*,&N^7T^3&\le P_*^3.
\end{align*}
\end{lemma}
\begin{proof}
All powers have positive bases. $N\le T\le N^2$ follows from
$1\le6/5$ and $5/3\le2$. From $T\ge N^{6/5}$,
$T^{5/7}\ge N^{6/7}$, giving $N\le L_*$.
From $T\le N^{5/3}$, $T^{3/35}\le N^{1/7}$, giving
$T^{4/5}=T^{3/35}T^{5/7}\le L_*$. Using $N\le T$,
$L_*\le T^{6/7}\le T$. Multiplying by five and
$5^{4/5}\le5$ yields $(5T)^{4/5}\le5L_*\le5T$.
Also $U_*\ge1$ and, using $T\le N^2$,
$U_*\le N^{6/7}\le N$. Every identity is an exact addition of
exponents; for example $NL_*U_*^2$ has $N$-exponent
$1+1/7+10/7=18/7$ and $T$-exponent $5/7+1/7=6/7$.
For the final inequality,
$N^7T^3/P_*^3=N^{-5/7}T^{3/7}\le1$ because $T^3\le N^5$.
\end{proof}
\source{GuthMaynard.Improvement.MiddleRangeScales; MiddleRangeIdentities; MiddleLargeValueBranches.middle_threshold_le_length}

\begin{lemma}[Homogeneous cubic absorption]\label{lem:absorb}
If $M,P,Q\ge0$ and
\[
 M^3\le Q(P^3+P^{3/2}M^{3/2}+P^2M),
\]
then $M\le(1+3Q)P$. If the right-hand coefficient is $QN^\eps$ with
$N\ge1$, $\eps\ge0$, then
$M\le(1+3Q)N^\eps P$.
\end{lemma}
\begin{proof}
If $M\le P$, the first conclusion is immediate. Otherwise $M>0$ and
$P\le M$. Each term in the bracket is at most $PM^2$:
$P^3\le PM^2$, $P^{3/2}M^{3/2}\le PM^2$, and $P^2M\le PM^2$.
Thus $M^3\le3QPM^2$; cancellation of $M^2>0$ gives $M\le3QP$.
For the second assertion use the first with $QN^\eps$ and then
$1+3QN^\eps\le(1+3Q)N^\eps$. The argument does not divide by $P$,
so it includes $P=0$.
\end{proof}
\source{GuthMaynard.Improvement.MomentScalarClosure.normalized_moment_absorption, normalized_moment_power_absorption}

\begin{theorem}[Actual optimized truncated fourth moment]\label{thm:fourth}
For every $\eps>0$ there is $C_\eps>0$ such that, uniformly for
$N\ge2$, $N^{6/5}\le T\le N^{5/3}$, and bounded coefficients,
\[
 \int_\R\left(\one_{\{|D_{N,a}|\le U_*\}}
                  \chi(t/T)|D_{N,a}(t)|^2\right)^2dt
 \le C_\eps N^\eps P_*.
\]
\end{theorem}
\begin{proof}
Write $L=L_*$, $U=U_*$, $P=P_*$ and $\theta=\eps/6$.
Use \eqref{eq:fixedraw} with $\delta=\theta$, cubic loss $\theta$,
and spectral scale $5L$. Lemma~\ref{lem:scales} turns it exactly into
\begin{align}\label{eq:middleraw}
 M^3\le SN^\theta\bigg[&125P^3+25P^2M+C(5T)^\theta\bigg(
        5P^{3/2}M^{3/2}\\
 &+5D_1N^4T(N^2+H_{\rm rem})M
               +3060D_1N^7T^3\log^2(2N)\bigg)\bigg].\nonumber
\end{align}
The $125$ and $25$ are the literal factor-five rescaling losses.
Use Proposition~\ref{prop:removed} with loss $\theta$, giving
\[
 H_{\rm rem}\le KN^\theta(N^2+NL)\le2KN^\theta NL,
 \qquad N^2+H_{\rm rem}\le(1+2K)N^\theta NL.
\]
The latter and $N^5LT=P^2$ bound the middle removal term by
$5D_1(1+2K)N^\theta P^2M$.
Moreover
\[
 (5T)^\theta\le5^\theta N^{2\theta},\qquad
 \log^2(2N)\le J^2N^{2\theta},\quad J=\log2+1/\theta,
 \qquad N^7T^3\le P^3.
\]
Since $N^\theta,N^{2\theta}\ge1$, all nonnegative terms in
\eqref{eq:middleraw} are dominated by
\[
 QN^{\theta+\theta+2\theta+2\theta}
       (P^3+P^{3/2}M^{3/2}+P^2M),
\]
where the source's explicit uniform coefficient is
\[
 Q=S\left(150+C5^\theta
       \left(5+5D_1(1+2K)+3060D_1J^2\right)\right)\ge0.
\]
The four losses sum to $6\theta=\eps$. Lemma~\ref{lem:absorb} gives
$M\le(1+3Q)N^\eps P$. Every constant was selected before $N,T,a$.
\end{proof}
\source{GuthMaynard.Improvement.MiddleMomentRaw.middle_source_moment_raw; MomentScalarClosure.raw_moment_normalized_budget; MiddleFourthMoment.middle_fourth_moment}

\chapter{From the fourth moment to all separated large values}\label{ch:middle}
\section{Good cells, bad cells, and the reproducing weight}
The source fixes the exact frequency cutoff $\chi$ of Chapter~\ref{ch:cubic}
and defines $k=\mathcal F^{-1}\chi$ as a Schwartz function. Put $q(t)=|k(t)|$;
no smoothness of the absolute value $q$ is required.
For $n\in I_N$ the negative frequency
$-\log(n/N)/(2\pi)$ lies in $[-1,2]$, so the transform of $k$ equals one
there. Translation invariance gives, term by term,
\[
 \int k(t-w)e(t\xi)\,dt=e(w\xi)\widehat k(-\xi).
\]
Finite summation therefore proves $D(w)=\int k(t-w)D(t)\,dt$.
This applies inversion to the Schwartz kernel, not to the nonintegrable
exponential. Two weighted Cauchy--Schwarz inequalities give
\[
 |D(w)|^4\le\left(\int q\right)^3\int q(t-w)|D(t)|^4dt.
\]
Schwartz decay supplies $q(s)\le C_j(1+|s|)^{-j}$. Applying the source's
separated quadratic-decay packing lemma gives the uniform overlap bound
$\sum_{w\in W}q(t-w)\le B_0$. Since $\int|s|^jq(s)\,ds$ is finite,
$q(s)\le |s|^jq(s)/R^j$ on $|s|>R$. Multiply by
$|D|^4\le16N^4$ and integrate after translating. This proves, for $R>0$,
\[
 \int_{|t-w|>R}q(t-w)|D(t)|^4dt\le D_jN^4/R^j.
\]
All weighted moments are integrable because $D$ is bounded and the kernel
is integrable; no global $L^4$ assertion for $D$ is used.
\source{GuthMaynard.Improvement.SourceReproducingKernel; SourceQuarticKernel; SourceQuarticTail}

Let
\[
 \mathcal K=\{k\in\Z:\lfloor-T\rfloor\le k\le\lfloor2T\rfloor,
       \ \exists x\in[-T,2T],\ |D(x)|>U,\ \lfloor x\rfloor=k\}.
\]
The exact deleted set is
$Z=\bigcup_{k\in\mathcal K}\{w\in W:|w-k|\le R+1\}$, and $G=W\setminus Z$.
A one-separated set in $[k-R-1,k+R+1]$ has at most $2R+3$ points.
Taking the union bound proves $|Z|\le(2R+3)|\mathcal K|$.
If $w\in G$ and $|t-w|\le R\le T$, then $t\in[-T,2T]$.
If $|D(t)|>U$, the integer $k=\lfloor t\rfloor$ belongs to $\mathcal K$,
and $|t-k|<1$ implies $|w-k|\le R+1$, contradicting $w\notin Z$.
Thus $|D(t)|\le U$ and $\chi(t/T)=1$, so
$\mu_U(t)^2=|D(t)|^4$. This explains the plateau $[-1,2]$.

For completeness, the bad-cell bound uses parity-bin selection alone.
Choose one witness from each occupied cell and select the parity class
containing at least half the cells. Those witnesses are one-separated.
Split them into the three intervals $[-T,0]$, $[0,T]$, $[T,2T]$;
each has length $T$, preserving $T\le N^2$. The classical bound applied
at threshold $U$ gives at most three copies of its budget. Restoring both
parities gives at most six copies. No derivative estimate or positive
lower bound on the measure of a bad cell is needed.
\source{GuthMaynard.Improvement.QuarticCellGeometry; QuarticMomentExtraction.quartic_good_integral}

\begin{proposition}[Extraction at a specified truncated moment]\label{prop:extraction}
For each $\eps>0$ there is $E_\eps>0$ such that if
$N\ge2$, $N\le T\le N^2$, $1\le V\le U\le2N$, and
$W\subset[0,T]$ is one-separated with $|D(w)|\ge V$, then for every
upper bound $M_0$ on $\int\mu_U^2$,
\begin{equation}\label{eq:extract}
 |W|\le E_\eps N^\eps
 \left(N^2/U^2+TN^4/U^6+M_0/V^4\right).
\end{equation}
\end{proposition}
\begin{proof}
For a fixed tail order $j$, the preceding reproducing estimates imply for
$w\in G$
\[
 V^4\le A_0\int q(t-w)\mu_U(t)^2dt+A_0D_jN^4/R^j,
 \qquad A_0=\left(\int q\right)^3.
\]
When the last term is at most $V^4/2$, sum over $G$, interchange the finite
sum and integrals, and use the overlap bound to obtain
$|G|\ll M_0/V^4$. The actual classical bad-cell bound gives
$|\mathcal K|\ll_{\eps}N^\eps(N^2/U^2+TN^4/U^6)$.
It uses occupied unit cells and the parity argument just proved, not the
measure of a superlevel set or an uncharged derivative-variation loss. Combining deleted and good points gives
\[
 |W|\le B\left((R+1)N^{\eps/2}(N^2/U^2+TN^4/U^6)+M_0/V^4\right)
\]
whenever its uniform rapid-tail condition holds.

Choose $\nu=\min(\eps/2,1/2)>0$ and an integer $j$ such that $\nu j\ge5$.
Let $N_0=\max(2,2B)$ and take $R=N^\nu$ for $N\ge N_0$.
Then $R\le N\le T$, $R^j\ge N^5$, and $BN^4/R^j\le1/2\le V^4/2$.
Also $(R+1)N^{\eps/2}\le2N^\eps$.
For $2\le N<N_0$, packing gives $|W|\le T+1\le2N_0^2$, while
$U\le2N$ gives $N^2/U^2\ge1/4$. Thus an enlarged common constant
$8N_0^2+2B+1$ covers the bounded-length case too. This proves
\eqref{eq:extract} with no radius or tail premise left.
\end{proof}
\source{GuthMaynard.Improvement.QuarticMomentExtraction.quartic_moment_extraction; QuarticBadCells.quartic_bad_cells_bound; QuarticExtraction.quartic_extraction; QuarticRadiusExtraction.quartic_extraction_of_moment}

\begin{lemma}[Threshold budgets]
In the middle time range, if $0<V\le U_*$, then
\[
 N^2/U_*^2+TN^4/U_*^6\le2P_*/V^4.
\]
If $V\ge U_*$, then $TN^4/V^6\le P_*/V^4$.
\end{lemma}
\begin{proof}
Since $N\le L_*$ and $P_*=NL_*U_*^2$,
$N^2/U_*^2\le P_*/U_*^4$.
The identity $TN^4/U_*^2=P_*$ gives
$TN^4/U_*^6=P_*/U_*^4$. Sum and use $V^4\le U_*^4$.
For the second assertion use $V^2\ge U_*^2$ in
$TN^4/V^6=(TN^4/V^2)/V^4$.
\end{proof}
\source{GuthMaynard.Improvement.MiddleLowLargeValues.middle_low_threshold_budget; MiddleLargeValueBranches.middle_high_threshold_budget}

\begin{proof}[Proof of Theorem~\ref{thm:middle}]
First treat $W\subset[0,T]$. If $0<V<1$, no polynomial estimate is needed:
packing gives $|W|\le T+1\le2N^2\le2N^2/V^2$.
If $1\le V\le U_*$, apply Proposition~\ref{prop:extraction} with loss
$\eps/2$ and Theorem~\ref{thm:fourth} with loss $\eps/2$.
The threshold-budget lemma gives
\[
 |W|\le E N^{\eps/2}
   \left(2P_*/V^4+D N^{\eps/2}P_*/V^4\right)
 \le E(2+D)N^\eps P_*/V^4.
\]
If $V>U_*$, apply the actual classical estimate
\eqref{eq:classical} with loss $\eps$ and the high-threshold budget.
This gives $|W|\le C_2N^\eps(N^2/V^2+P_*/V^4)$.
The sum $C=C_1+C_2+2$ of the low, high, and packing constants works for
all positive thresholds.

For a translated interval set
$W_0=\{t-t_0:t\in W\}$ and
$b_n=a_n\e^{\ii t_0\log(n/N)}$.
Translation preserves cardinality and separation, $|b_n|=|a_n|$, and
\[
 D_{N,b}(s)=D_{N,a}(s+t_0).
\]
Apply the origin bound to $W_0,b$. The right side and constant are unchanged.
\end{proof}
\source{GuthMaynard.Improvement.MiddleLowLargeValues.middle_large_values_low; MiddleLargeValueBranches.middle_large_values_high, middle_large_values_small; SourceTranslation; MiddleLargeValues.middle_range_large_values}

\chapter{Short zero detection and fixed coefficient families}\label{ch:detector}
\section{The arithmetic mollifier and the exact contour identity}
Let $\mu_{\rm ar}(n)$ be the integer M\"obius function, extended by zero at
$n=0$. The mollifier and its product coefficients are
\[
 M_X(s)=\sum_{1\le d\le\lfloor X\rfloor_+}\mu_{\rm ar}(d)d^{-s},\qquad
 c_X(n)=\sum_{d\mid n,\ d\le\lfloor X\rfloor_+}\mu_{\rm ar}(d).
\]
The divisor set at $0$ is empty. For $X\ge1$, $c_X(1)=1$ and
$c_X(n)=0$ if $2\le n\le X$. Also $|c_X(n)|\le d(n)$.
Define, for $Y>0$,
\[
 S(\rho;X,Y)=\sum_{n\ge0}c_X(n)n^{-\rho}\e^{-n/Y},\qquad
 R(\rho;X,Y)=M_X(1)Y^{1-\rho}\Gamma(1-\rho),
\]
where the $n=0$ summand is zero. The source proves absolute convergence
in the real-part ranges in which it is used. With $\rho=\beta+\ii\gamma$,
put
\begin{align*}
 I(\rho;X,Y)=\frac1{2\pi}\int_\R&
 \Gamma(1/2-\beta+\ii v)\zeta(1/2+\ii(\gamma+v))\\
 &\cdot M_X(1/2+\ii(\gamma+v))Y^{1/2-\beta+\ii v}\,dv.
\end{align*}
\source{MathCollab.Density.MollifierCoefficients; DetectorMellin.detectorKernel; DetectorContour.detectorResidue; DetectorError.detectorIntegral}

\begin{analyticinput}[Actual short contour identity]\label{input:contour}
If $\zeta(\rho)=0$, $9/13\le\beta<1$, and $Y\ge1$, then, for independent
mollifier and smoothing parameters $X,Y$,
\[
 S(\rho;X,Y)=R(\rho;X,Y)+I(\rho;X,Y).
\]
\end{analyticinput}
\paragraph{What the contour argument uses.}
Mellin inversion of $\e^{-n/Y}$ gives the right vertical contour on
$\Rea s=1/2$. Move to $\Rea s=1/2-\beta$ in a finite rectangle.
The factor $\zeta(\rho)=0$ cancels the Gamma singularity at $s=0$;
the zeta pole at $s=1-\rho$ gives exactly $R$.
The shifted zeta regularization, residue identity, absolute integrability
of both vertical contours, and decay of both horizontal contours are
proved dependencies. Taking the height to infinity gives the displayed
identity with factor $1/(2\pi)$, not $1/(2\pi\ii)$ after real vertical
parameterization. No Weyl estimate is an extra premise of this identity.
\source{GuthMaynard.Improvement.ShortDetectorIdentity.short_smoothedDetector_contour_identity, short_zeta_detector_identity; ShortDetectorContour; ShortDetectorMellin; ShortDetectorEstimates}

\section{The finite detector and its dichotomy}
For a positive height $H$ choose
\[
 X_H=2H^{1/100},\quad Y_H=H^{1/2},\quad
 K_H=\left\lceil H^{1/2}(\log H)^2\right\rceil_+.
\]
Define fixed coefficients
\[
 b_H(n)=\begin{cases}c_{X_H}(n)\e^{-n/Y_H},&
              \lfloor X_H\rfloor_+<n<K_H,\\0,&\text{otherwise},\end{cases}
 \qquad F_H(\rho)=\sum_{n<K_H}b_H(n)n^{-\rho}.
\]
They depend on $H$, not on either coordinate of the zero.
For $H\ge1$, $K_H>1$, and $\beta\ge0$, the source proves the exact split
\begin{equation}\label{eq:finiteidentity}
 S(\rho;X_H,Y_H)=\e^{-1/Y_H}+F_H(\rho)
          +\sum_{n\ge0}c_{X_H}(n+K_H)(n+K_H)^{-\rho}\e^{-(n+K_H)/Y_H}.
\end{equation}
The finite source interval is strict at both $X_H$ and
$Y_H(\log H)^2$, because the upper index is $n<K_H$.
\source{GuthMaynard.Improvement.ShortFiniteDetector: shortDetectorCoefficient, shortFiniteDetector, shortFiniteDetector_decomposition}

\begin{analyticinput}[Uniform short-detector remainder bounds]\label{input:remdet}
For every fixed $c>0$, eventually in $H$, the infinite tail in
\eqref{eq:finiteidentity} has norm at most $c$ uniformly for $\beta\ge0$;
the residue $R(\rho;X_H,Y_H)$ has norm at most $c$ uniformly for
$9/13\le\beta\le1$ and $H\le|\gamma|\le2H$.
\end{analyticinput}
These estimates use the exponential smoothing tail and actual Gamma-strip
bounds. In particular, the negative-strip bound keeps its singular real-part
factor explicit:
\[
 |\Gamma(-a+\ii v)|\le(12/a)(1+|v|)\e^{-|v|}
 \quad(0<a\le1/2).
\]
Here $a=\beta-1/2\ge5/26$ for the short branch. The source obtains this
negative-strip bound from $\Gamma(z+1)=z\Gamma(z)$ and $|z|\ge a$.
\source{GuthMaynard.Improvement.ShortDetectorTruncation.short_detector_tail_eventually_small; ShortDetectorResidue.short_detector_residue_eventually_small; ShortDetectorGamma.short_detector_gamma_bound}

\begin{lemma}[Uniform detector alternative]\label{lem:alternative}
Eventually in $H$, every zero $\rho$ with
$9/13\le\Rea\rho<1$, $H\le|\Ima\rho|\le2H$ satisfies
\[
 |F_H(\rho)|\ge1/3\quad\text{or}\quad
 |I(\rho;X_H,Y_H)|\ge1/3.
\]
\end{lemma}
\begin{proof}
Take both remainders in Input~\ref{input:remdet} to be at most $1/24$,
and $Y_H\ge4$. Since $\e^{-1/Y_H}\ge1-1/Y_H\ge3/4$, the contour
identity and \eqref{eq:finiteidentity} give
$3/4\le|F_H(\rho)|+|I(\rho;X_H,Y_H)|+1/12$.
If both displayed detector values were less than $1/3$, the right side
would be less than $2/3+1/12=3/4$, a contradiction.
All required thresholds precede $\rho$.
\end{proof}
\source{GuthMaynard.Improvement.ShortDetectorAlternative.short_detector_eventually_alternative}

\section{Eliminating the varying real part by an exact Taylor family}
Put $D_H=\lfloor\log_2 K_H\rfloor+1$ in the natural-logarithm-to-base-two
integer convention \lean{Nat.log}. Decompose the finite detector into
blocks $(Q,2Q]$, $Q=2^r$, $0\le r<D_H$:
\[
 B_{H,Q}(\rho)=\sum_{Q<n\le2Q}b_H(n)n^{-\rho}.
\]
If $|F_H(\rho)|\ge1/3$, a block has norm at least $1/(3D_H)$.
Every nonzero block has
$H^{1/100}<Q<H^{1/2}(\log H)^2$.
For $j\ge0$ define
\[
 b_{H,Q,j}(n)=b_H(n)(2n/Q-3)^j,\qquad
 P_{H,Q,j}(t)=\sum_{Q<n\le2Q}b_{H,Q,j}(n)\e^{-\ii t\log(n/Q)}.
\]
Again the coefficients do not depend on $\beta$ or $\gamma$.
Let
\[
 w_j(\beta)=\frac{\beta(\beta+1)\cdots(\beta+j-1)}{j!}(-1/3)^j,
 \qquad w_0(\beta)=1.
\]
For $0\le\beta\le1$ the rising-product ratio lies in $[0,1]$, so
$|w_j(\beta)|\le3^{-j}$. The binomial series gives the exact identity
\[
 \sum_{j\ge0}w_j(\beta)P_{H,Q,j}(\gamma)
       =(3/2)^\beta\sum_{Q<n\le2Q}b_H(n)(n/Q)^{-\beta}
                                    \e^{-\ii\gamma\log(n/Q)}.
\]
Indeed $1+(2n/Q-3)/3=2n/(3Q)$.
The norm of the right side is $(3/2)^\beta Q^\beta|B_{H,Q}(\rho)|$.
\source{MathCollab.Density.DetectorTaylorFamily.detectorTaylorWeight; GuthMaynard.Improvement.ShortDetectorTaylorFamily; ShortDetectorBlockExpansion.hasSum_shortDetectorTaylorPolynomial}

\begin{proposition}[A fixed logarithmic-size detecting family]\label{prop:fixedfamily}
Let $J_H=\lceil10\log H\rceil_+$. Eventually in $H$, the fixed family
$(r,j)$, $0\le r<D_H$, $0\le j\le J_H$, has size at most
$60(\log H)^2$. If $0\le\sigma\le\beta\le1$ and
$|F_H(\rho)|\ge1/3$, one member with $Q=2^r$ satisfies
\[
 |P_{H,Q,j}(\gamma)|\ge Q^\sigma/(60\log H),\qquad
 H^{1/100}<Q<H^{1/2}(\log H)^2.
\]
\end{proposition}
\begin{proof}
On $(Q,2Q]$, $|2n/Q-3|\le1$, $|b_H(n)|\le d(n)\le n\le2Q$;
therefore $|P_{H,Q,j}(t)|\le2Q^2$ for every $j,t$.
The infinite Taylor tail after $J$ has norm at most $3Q^2\,3^{-J}$.
At $J=J_H+1$, $3^{-J}\le H^{-5}$, and $Q\le H^2$ makes it at most
$3/H$. Eventually $3/H\le1/(6D_H)$.
For the block of norm $\ge1/(3D_H)$ the full Taylor combination has
norm at least $Q^\beta/(3D_H)$. Removing the tail leaves at least
$Q^\beta/(6D_H)$. The total absolute Taylor weight is at most
$\sum3^{-j}=3/2\le2$, so one retained polynomial has norm at least
$Q^\beta/(12D_H)$. Use $Q^\sigma\le Q^\beta$ and
$D_H\le5\log H$. Finally $J_H+1\le12\log H$, giving the stated family
cardinality. No coefficient is chosen after seeing the zero.
\end{proof}
\source{GuthMaynard.Improvement.ShortFixedDetectorFamily.shortFiniteDetector_fixed_family_eventually; ShortDetectorFamilyScales.short_detector_family_scales_eventually}

\section{Powering, dyadic product pieces, and normalization}
Extend $b_{H,Q,j}$ by zero outside $(Q,2Q]$ to an arithmetic function.
Its $k$-fold Dirichlet convolution has coefficients
\[
 c_{H,Q,j,k}(n)=\sum_{n_1\cdots n_k=n}\prod_{i=1}^k b_{H,Q,j}(n_i).
\]
It vanishes unless $Q^k<n\le(2Q)^k$ for $k\ge1$.
The finite Dirichlet series is exactly the $k$th power of the original
series. Use conjugated coefficients $\overline{c_{H,Q,j,k}(n)}$ so that
positive phases give the same norm at the \emph{same} ordinate $\gamma$.
Splitting $(Q^k,2^kQ^k]$ into $k$ dyadic intervals with starts
$N=2^\ell Q^k$, $0\le\ell<k$, one interval has polynomial norm at least
$|P_{H,Q,j}(\gamma)|^k/k$.
\source{GuthMaynard.Improvement.ShortDetectorConvolution; ShortDetectorPoweredPieces.exists_shortPoweredDetector_piece}

\begin{lemma}[Least-successful powering]
Uniformly for $a\in[3/4,4/5]$ and large $H$, if
$H^{1/100}<Q\le H^{1/2}(\log H)^2$, there is $1\le k\le100$ with
$H^a\le Q^k\le H^{3a/2}$. Every product piece satisfies
\[
 H^a\le N\le2^{99}H^{3a/2},\qquad2N\le H^2.
\]
\end{lemma}
\begin{proof}
Since $Q^{100}>H\ge H^a$, choose the least successful $k$. If $k\le2$,
$Q^k\le Q^2\le H(\log H)^4\le H^{3a/2}$ eventually, uniformly because
$3a/2\ge9/8>1$. If $k\ge3$, leastness gives $Q^{k-1}<H^a$;
$k/(k-1)\le3/2$ gives $Q^k<H^{3a/2}$. For the dyadic pieces use
$1\le2^\ell\le2^{99}$. The uniform buffer
$2^{100}H^{3a/2}\le H^2$ follows from $3a/2\le6/5<2$.
\end{proof}
\source{GuthMaynard.Improvement.ShortDetectorPowerScales.short_detector_power_exists, short_detector_power_buffer_eventually, short_detector_power_piece_bounds}

For fixed $k\le100$, divisor growth gives
$|c_{H,Q,j,k}(n)|\ll_\alpha n^\alpha$, uniformly in $H,Q,j,k$.
For any $\eta>0$, when $n\le H^2$, eventually these coefficients are
at most $H^\eta$. Define the normalized coefficients
\[
 a_{H,Q,j,k,\eta}(n)=\overline{c_{H,Q,j,k}(n)}/H^\eta.
\]
They have modulus at most one on all eligible pieces.
The full index set is
\[
 \mathcal P_H=\{(r,j,k,\ell):0\le r<D_H,\ 0\le j\le J_H,
                                  1\le k\le100,\ 0\le\ell<100\}.
\]
Its scale is $N_p=2^\ell(2^r)^k$ and its coefficient is the preceding
normalized function. Including extra pieces with $\ell\ge k$ in this fixed
index set is harmless; the actual covering witness has $\ell<k$.

\begin{proposition}[The powered fixed-family cover]\label{prop:poweredcover}
For each $\eta>0$, eventually in $H$, $|\mathcal P_H|\le600000(\log H)^2$.
For $a\in[3/4,4/5]$, $0\le\sigma\le\Rea\rho\le1$, and
$|F_H(\rho)|\ge1/3$, there is $p\in\mathcal P_H$ such that
\[
 H^a\le N_p\le2^{99}H^{3a/2},\qquad2N_p\le H^2,
 \qquad |D_{N_p,a_{p,\eta,H}}(\Ima\rho)|\ge N_p^\sigma H^{-2\eta}.
\]
\end{proposition}
\begin{proof}
Use Proposition~\ref{prop:fixedfamily}, the least-successful power, and the
piece pigeonhole. The factors $(60\log H)^k$, $k$, and the fixed conversion
from $Q^{k\sigma}$ to $(2^\ell Q^k)^\sigma$ are bounded by a constant
times $(\log H)^{100}$ uniformly for $k\le100$, $0\le\sigma\le1$.
For each fixed $\eta>0$, the source proves that this entire height loss is
at most $H^\eta$ eventually. Division by the coefficient normalization
$H^\eta$ gives $N_p^\sigma H^{-2\eta}$. The family cardinality is at most
$100\cdot100$ times $60(\log H)^2$. The coefficient bound is uniform over
all family members with $2N_p\le H^2$. This construction preserves the
ordinate, avoiding an uncharged reseparation loss.
\end{proof}
\source{GuthMaynard.Improvement.ShortPoweredDetectorFamily.shortPoweredDetectorFamily_cover, shortPoweredDetectorFamily_coefficients; ShortDetectorPoweredHeight; ShortDetectorNormalization}

\chapter{The two density windows from the improved large-values estimate}\label{ch:windows}
\section{The variable short-window exponents}
For $9/13\le\sigma\le5/7$, define
\[
 a(\sigma)=\frac6{7\sigma+3},\qquad b(\sigma)=\frac32a(\sigma),\qquad
 g(\sigma)=\frac{18(1-\sigma)}{7\sigma+3}.
\]
The source proves
\[
 3/4\le a\le4/5,\quad0\le b\le6/5,\quad g>0,
 \qquad2b(1-\sigma)=g,\qquad a(18/7-4\sigma)+6/7=g.
\]
Also
\begin{equation}\label{eq:meancompare}
 11/6-5\sigma/3\le g(\sigma).
\end{equation}
For completeness, multiplying the difference in \eqref{eq:meancompare}
by $6(7\sigma+3)>0$ gives
$70\sigma^2-155\sigma+75
=70(\sigma-5/7)(\sigma-3/2)\ge0$.
\source{GuthMaynard.Improvement.ShortDensityScale.short_density_scale_bounds, short_density_scale_identities}

\begin{proposition}[Large values in the variable powered window]\label{prop:shortwindow}
For $0<\delta\le1/100$ there is a positive constant independent of $\sigma$
such that, for $H,N\ge2$, $9/13\le\sigma\le5/7$,
\[
 H^{a(\sigma)}\le N\le H^{b(\sigma)+\delta},\qquad
 V=N^\sigma H^{-2\delta},
\]
and any bounded source, one-separated $W\subset[t_0,t_0+H]$ with
$|D(w)|\ge V$ obeys
\[
 |W|\le C_\delta H^{g(\sigma)+10\delta}.
\]
\end{proposition}
\begin{proof}
We have $N^\delta\le H^{2\delta}$ and
$N^{2-2\sigma}\le H^{g+2\delta}$.
If $N\le H^{5/6}$, then $N^{6/5}\le H$.
The lower bound $N\ge H^a\ge H^{3/5}$ gives $H\le N^{5/3}$.
Thus Theorem~\ref{thm:middle} is used in its exact time range. Its two
terms, after division by the threshold, are
\[
 N^\delta\left(N^{2-2\sigma}H^{4\delta}
             +N^{18/7-4\sigma}H^{6/7+8\delta}\right).
\]
Since $18/7-4\sigma\le0$, the lower length bound and the scale identity
give $N^{18/7-4\sigma}H^{6/7}\le H^g$.
The first term costs at most $H^{g+6\delta}$ inside the parentheses;
the second at most $H^{g+8\delta}$. The outer $N^\delta$ costs
$H^{2\delta}$, proving this branch with $2C H^{g+10\delta}$.

If $N>H^{5/6}$, use the actual separated second moment. Division by
$V^2$ gives
\[
 |W|\le102\left(HN^{1-2\sigma}+N^{2-2\sigma}\log(2N)\right)H^{4\delta}.
\]
Because $1-2\sigma\le0$,
$HN^{1-2\sigma}\le H^{11/6-5\sigma/3}\le H^g$.
With $D_\delta=\log2+1/\delta$, the logarithm is at most
$D_\delta N^\delta\le D_\delta H^{2\delta}$, so the second term costs
at most $D_\delta H^{g+4\delta}$. This branch costs at most
$102(1+D_\delta)H^{g+8\delta}$ and hence the claimed bound.
The source chooses $2C+102(1+D_\delta)$ as a common constant.
\end{proof}
\source{GuthMaynard.Improvement.ShortWindowLargeValues.short_window_large_values; ShortDensityHeight.short_density_height_budgets}

\section{The closed upper window and the padded time}
For $5/7\le\sigma\le3/4$ set
\[
 h(\sigma)=\frac94(1-\sigma),\qquad
 s(\sigma)=\begin{cases}6/7,&\sigma\le18/25,\\7/8,&\sigma>18/25.\end{cases}
\]
The exact scalar comparisons are
\[
 39/14-3\sigma\le h,\quad
 s(18/5-4\sigma)\le h,\quad1+s(1-2\sigma)\le h.
\]
The first is equality at $5/7$. For the $6/7$ split the respective latter
differences are $(-117+165\sigma)/140$ and $(11-15\sigma)/28$;
for the $7/8$ split they are $(-18+25\sigma)/20$ and $(3-4\sigma)/8$.
They are nonnegative on the specified subintervals.
\source{GuthMaynard.Improvement.UpperCurveBudgets.upper_curve_exponents}

\begin{proposition}[Large values in the closed upper powered window]\label{prop:upperwindow}
For $0<\delta\le1/100$ there is one constant for all
$5/7\le\sigma\le3/4$, $H,N\ge2$, and
$H^{3/4}\le N\le H^{9/8+\delta}$, with the same bounded coefficients,
separation, localization, and threshold $V=N^\sigma H^{-2\delta}$:
\[
 |W|\le C_\delta H^{h(\sigma)+10\delta}.
\]
\end{proposition}
\begin{proof}
As before $N^\delta\le H^{2\delta}$ and
$N^{2-2\sigma}\le H^{h+2\delta}$.
If $N\le H^{s(\sigma)}$, use the padded time
\[
 S=\begin{cases}H,&N\le H^{5/6},\\N^{6/5},&N>H^{5/6}.\end{cases}
\]
It satisfies $H\le S$ and $N^{6/5}\le S\le N^{5/3}$.
When $S=H$, the lower bound $N\ge H^{3/4}$ gives
$N^{18/7-4\sigma}S^{6/7}\le H^{39/14-3\sigma}\le H^h$.
When $S=N^{6/5}$, the term is $N^{18/5-4\sigma}$; its exponent is
nonnegative, so the upper length bound gives at most
$H^{s(18/5-4\sigma)}\le H^h$.
Apply Theorem~\ref{thm:middle} on the enlarged interval
$[t_0,t_0+S]$. The original set and threshold are unchanged.
The same two-term calculation as in Proposition~\ref{prop:shortwindow}
gives at most $2C H^{h+10\delta}$.
For $N>H^{s(\sigma)}$, the separated second moment and the last scalar
comparison give $HN^{1-2\sigma}\le H^h$; the logarithmic and first-term
budgets give the same mean-value closure. This proof includes $\sigma=3/4$.
\end{proof}
\source{GuthMaynard.Improvement.UpperCurveBudgets.upper_curve_padded_time, upper_curve_mean_budget; UpperCurveLargeValues.upper_curve_large_values}

\section{Type-I separated zero representatives}
\begin{proposition}[Type-I counts in both windows]\label{prop:typeI}
For each $\eps>0$, uniformly in $\sigma,t_0$ and eventually in $H$,
any finite $R\subset\C$ whose ordinates are injective and one-separated,
with $\sigma\le\Rea\rho\le1$, $\Ima\rho\in[t_0,t_0+H]$ and
$|F_H(\rho)|\ge1/3$, has
\[
 |R|\le C_\eps H^{g(\sigma)+\eps}\quad(9/13\le\sigma\le5/7),
 \qquad |R|\le C_\eps H^{h(\sigma)+\eps}\quad(5/7\le\sigma\le3/4).
\]
\end{proposition}
\begin{proof}
Take $\eta=\min(\eps/100,1/100)$. For the short window apply
Proposition~\ref{prop:poweredcover} with $a=a(\sigma)$;
for the upper window use $a=3/4$.
Restrict the fixed family to members satisfying the appropriate length
bounds and $2N_p\le H^2$. Eventually $2^{99}\le H^\eta$, so its upper
length becomes $H^{b+\eta}$ or $H^{9/8+\eta}$ respectively.
For each member let $W_p$ be the subset of ordinates where that member
has norm at least $N_p^\sigma H^{-2\eta}$.
The corresponding window proposition bounds
$|W_p|\le C_\eta H^{g+10\eta}$ or $C_\eta H^{h+10\eta}$.
Each $W_p$ is a subset of the original separated set, so separation and
localization are preserved. The cover gives $W\subset\bigcup_pW_p$.
Eventually $600000(\log H)^2\le H^\eta$. Hence
$|R|=|W|\le\sum_p|W_p|\le C_\eta H^{g+11\eta}$, or the version with
$h$. Since $11\eta\le\eps$, the desired estimate follows.
All finite families and their coefficients were constructed before any zero
or real part was chosen.
\end{proof}
\source{GuthMaynard.Improvement.ShortTypeICount.short_typeI_separated_count, short_typeI_zero_slab_count; UpperCurveTypeICount.upper_curve_typeI_separated_count, upper_curve_typeI_zero_slab_count}

\section{The actually used Type-II moment is weaker than the named contract}
Put
\[
 Z_H(t)=\left|\zeta(1/2+\ii t)M_{2H^{1/100}}(1/2+\ii t)\right|.
\]
The proposition \lean{ShortMobiusFourthMoment} is the assertion
$\int_{H/2}^{3H}Z_H^4\ll_\eps H^{1+\eps}$.
Some conditional intermediate modules explicitly take that assertion as a
hypothesis. \emph{The final short and upper-curve branches do not use it.}
They instead use the following proved weaker theorem.

\begin{analyticinput}[Native ordinary fourth moment]\label{input:fourthzeta}
For every $\eps>0$ there are $B>0$ and a height threshold such that
\[
 \int_x^{2x}|\zeta(1/2+\ii t)|^4dt\le Bx^{1+\eps}
\]
for all positive $x$ above that threshold.
\end{analyticinput}
This is the native actual-zeta producer
\lean{Fourth.zeta_fourth_dyadic}, not a postulated moment inequality.
The following expands the actual producer above its already established
contour and Gamma-ratio machinery.
\source{MathCollab.Density.Stronger.Fourth.FourthMoment.zeta_fourth_dyadic}

\paragraph{Exact divisor-square source.}
Let $s_t=1/2+\ii t$, $d(n)$ be the divisor count ($d(0)=0$),
$\Gamma_\R(s)=\pi^{-s/2}\Gamma(s/2)$,
$G_t=\Gamma_\R(s_t)\Gamma_\R(s_{-t})\ne0$, and
$A(w)=(1-4w^2)^4$. Define
\[
 P_t(w)=\left(\frac{(s_{-t}+w)(1-s_{-t}-w)}{s_{-t}(1-s_{-t})}\right)^2,
\]
\[
 K_t(w)=\frac{\e^{100w^2}A(w)P_t(w)\Gamma_\R(s_{-t}+w)^2}{G_tw}.
\]
For $c>0$ put
\[
 C_{t,c}(n)=\frac1{2\pi}\int_\R d(n)n^{-s_{-t}-c-\ii u}K_t(c+\ii u)\,du,
 \qquad P_{t,c,N}=\sum_{0\le n\le N}C_{t,c}(n),
\]
with the zero-index term explicitly zero. The exact source's right piece
$R_t$ is its vertical divisor integral at real coordinate $1$ divided by
$G_t$. The native contour reflection and completed-zeta identity give
$z(t)^2=2\Rea R_t$, hence $z(t)^4\le4|R_t|^2$.
For $t\ge0$, each $C_{t,c}(n)$ is independent of $c>0$: a finite rectangle
has zero boundary integral, both vertical lines are integrable, and both
horizontal integrals tend to zero. The source proves
$R_t=\sum_{n\ge0}C_{t,c}(n)$ and the exact finite-prefix-plus-tail split.
These are the precise contour identities used here, not an assumed
approximate functional equation with unnamed errors.
\source{MathCollab.Density.Stronger.Fourth.SquareSource; SquareNormalized; SquareLineShift; SquareTruncation}

\paragraph{Truncation and Gaussian prefix control.}
For $b>3/2$ the established Gamma-kernel bound and $d(n)\le n$ give
\[
 |R_t-P_{t,c,N}|\le A_b t^b N^{3/2-b}
\]
when $t\ge\max(4,4b)$, for every $c>0$.
The tail sum is bounded by the decreasing-power integral
$\sum_{n>N}n^{1/2-b}\le N^{3/2-b}/(b-3/2)$.
For any $\delta>0$ choose $b=2+2/\delta$; then
$b+(1+\delta)(3/2-b)=-(1+\delta)/2\le0$.
If $H\le t\le2H$ and $N\ge H^{1+\delta}$, this proves a uniform constant
bound $|R_t-P_{t,c,N}|\le K_{\rm tail}$.

For fixed $c>0$ the kernel estimate is
$|K_t(c+\ii u)|\le A_ct^c\e^{-98u^2}$ for $t\ge\max(4,4c)$.
Write
$a_c(n)=d(n)n^{-1/2-c}$ and
$D_{c,N}(v)=\sum_{1\le n\le N}a_c(n)\e^{-\ii v\log n}$.
The exact finite-prefix integral and weighted Cauchy--Schwarz give
\[
 |P_{t,c,N}|^2\le K_g t^{2c}
       \int_\R\e^{-98u^2}|D_{c,N}(u-t)|^2du.
\]
Both weighted integrals are finite since $D_{c,N}$ is a bounded finite sum,
and the Gaussian mass is $\sqrt{\pi/98}$.
The divisor subpower estimate gives, uniformly in $c\ge0$,
$|a_c(n)|^2\le D_\theta n^{\theta-1}$; summing on $(M,2M]$ gives
$\sum|a_c(n)|^2\le D_\theta(2M)^\theta$.
The proved dyadic-prefix mean square, followed by Gaussian Fubini, therefore
supplies, for $N=2^M$,
\[
 \int_H^{2H}\int_\R\e^{-98u^2}|D_{c,N}(u-t)|^2du\,dt
 \le D_p(M+1)^2(H+2(5\pi+1)N)N^\theta.
\]
\source{MathCollab.Density.Stronger.Fourth.SquareTail; SquarePrefix; Coefficients; PrefixMean; GaussianProducer}

\paragraph{The actual fourth-moment exponent closure.}
For requested loss $\eps>0$ the source's expression
\lean{min eps 1/20} parses as $q=\min(\eps,1)/20$.
In particular $0<q\le1/20$, $q\le1$, and $7q\le\eps$.
Choose $H^{1+q}\le N=2^M\le2H^{1+q}$.
The tail and prefix estimates give
\[
 z(t)^4\le8K_g(2H)^{2q}
          \int_\R\e^{-98u^2}|D_{q,N}(u-t)|^2du+8K_{\rm tail}^2.
\]
The dyadic-count estimate $(M+1)^2\le A_qN^q$ and
$H+2(5\pi+1)N\le(1+4(5\pi+1))H^{1+q}$ show the integrated main budget is
at most a constant times
\[
 (2H)^{2q}H^{1+q}(2H^{1+q})^{2q}
       =2^{4q}H^{1+5q+2q^2}\le2^{4q}H^{1+7q}.
\]
The integrated constant tail is $8K_{\rm tail}^2H$.
Since $7q\le\eps$ and $H\ge1$, both fit $CH^{1+\eps}$.
The height threshold is fixed before $H$; all finite-prefix, Gaussian, and
coefficient constants were already fixed by $q$. This proves exactly
Input~\ref{input:fourthzeta}, relative to the explicit contour-kernel inputs.
\source{MathCollab.Density.Stronger.Fourth.FourthMoment.zeta_fourth_dyadic; GaussianProducer.exists_fourth_dyadic_budget}

\begin{proposition}[Weak short mollified fourth moment]\label{prop:weakmoment}
For every $\eps>0$, eventually in $H$,
\[
 \int_{H/2}^{3H}Z_H(t)^4dt\le C_\eps H^{1+1/25+\eps}.
\]
\end{proposition}
\begin{proof}
The elementary mollifier bound at $\Rea s=1/2$ is
$|M_X(s)|\le\lfloor X\rfloor_+\le X$. Thus
$Z_H(t)^4\le16H^{1/25}|\zeta(1/2+\ii t)|^4$.
Cover $[H/2,3H]$ by $[H/2,H]$, $[H,2H]$, and $[2H,4H]$.
For sufficiently large $H$ the native fourth-moment bound applies on all
three intervals. Their sum is at most $3(B+1)(2H)^{1+\eps}$.
Multiplying by $16H^{1/25}$ gives the source's constant
$48(B+1)2^{1+\eps}$. Closed versus half-open integration endpoints have
Lebesgue measure zero.
\end{proof}
\source{GuthMaynard.Improvement.WeakShortMoment.shortTwistedNorm_fourth_le, weak_short_moment}

\begin{analyticinput}[Window reduction for actual Type-II detectors]\label{input:typeIIwindow}
There is $C>0$, eventually in $H$, uniformly for $\sigma\ge9/13$,
one-separated $W\subset[H,2H]$, and a chosen zero representative $\rho_u$
with $\Ima\rho_u=u$, $\sigma\le\Rea\rho_u\le1$, such that
$|I(\rho_u;X_H,Y_H)|\ge1/3$ implies
\[
 |W|\le CH^{1-2\sigma}\int_{H/2}^{3H}Z_H(t)^4dt.
\]
\end{analyticinput}
\paragraph{Precise reduction.}
The detector-window theorem, after Gamma-tail truncation and H\"older,
supplies for each $u$ a lower bound
\[
 1\le A H^{1-2\sigma}\int_{-H/2}^{H/2}
 |\Gamma(1/2-\Rea\rho_u+\ii v)|Z_H(u+v)^4dv.
\]
Sum over the one-separated ordinates. The uniform negative-Gamma strip
bound and the proved separated Gamma-packing bound dominate the sum of
weights by a constant on $[H/2,3H]$. This gives the displayed input.
Different representatives are allowed to have different real parts;
no common-$\beta$ assumption is added.
\source{GuthMaynard.Improvement.ShortDetectorWindowDetection.short_detector_window_fourth; ShortTypeIIPacking.short_typeII_window_packing; ShortTypeIICount.short_typeII_count_by_moment}

\begin{proposition}[The unconditional Type-II count used here]\label{prop:typeII}
For every $\eps>0$, eventually in $H$, Type-II separated representatives
in a positive slab have
\[
 |R|\le C_\eps H^{2-2\sigma+1/25+\eps}.
\]
On $9/13\le\sigma\le5/7$ this exponent without $\eps$ is at most
$g(\sigma)$; on $5/7\le\sigma\le3/4$ it is at most $h(\sigma)$.
\end{proposition}
\begin{proof}
Multiply Input~\ref{input:typeIIwindow} by Proposition~\ref{prop:weakmoment}.
The exponent is $(1-2\sigma)+(1+1/25+\eps)$.
For the first comparison multiply $g-(2-2\sigma+1/25)$ by
$25(7\sigma+3)>0$. The exact polynomial is
\[
 350\sigma^2-657\sigma+297
 =\frac{25}{14}(7\sigma-5)(28\sigma-33)
       +\frac{11}{2}\sigma+\frac{33}{14}.
\]
On $9/13\le\sigma\le5/7$, both $7\sigma-5$ and $28\sigma-33$ are
nonpositive, so their product is nonnegative; the remaining two terms are
positive. This is the nonpositive-factor witness actually used by
\lean{weak_short_typeII_exponent_le}, with the scalar algebra made explicit.
In particular, at $\sigma=5/7$ the scaled difference is $44/7$, not zero.
For the second comparison the difference is
$(1-\sigma)/4-1/25\ge1/16-1/25>0$ on the closed upper interval.
Thus the genuine fixed power loss $1/25$ fits both target windows.
\end{proof}
\source{GuthMaynard.Improvement.WeakTypeIICount.weak_short_typeII_count; WeakShortDensity.weak_short_typeII_zero_slab_count, weak_short_typeII_exponent_le; UpperCurveDensity.upper_curve_slab_density}

\chapter{Restoring analytic multiplicities and summing the height slabs}\label{ch:counting}
\section{Local Jensen counts and separated representatives}
\begin{analyticinput}[The precise local analytic bounds]\label{input:jensen}
For $\sigma\ge3/5$, $H\ge8$, and an occupied ordinate bin $[k,k+1)$,
center the disks at $c=2+\ii(k+1/2)$. All slab zeros in that bin are inside
the closed disk of radius $3/2$. Zeta is analytic on the disk of radius
$7/4$, has $|\zeta(c)|\ge0.6$, and on its outer circle
$|\zeta|\le100H^3$.
Jensen's divisor bound with analytic multiplicity gives
\[
 \sum_{\substack{\rho\in Z(\sigma;H,2H)\\k\le\Ima\rho<k+1}}
 m(\rho)
 \le\frac{\log(100H^3/0.6)}{\log((7/4)/(3/2))}.
\]
\end{analyticinput}
The geometric inclusion follows from
$|\Rea(\rho-c)|\le7/5$, $|\Ima(\rho-c)|\le1/2$, and
$(7/5)^2+(1/2)^2< (3/2)^2$.
The outer disk avoids the pole $1$ at these heights. The lower Euler-product
bound, outer zeta-growth bound, and analytic-order/divisor identification
are the exact imported analytic facts used by Jensen.
\source{GuthMaynard.Improvement.ShortLocalZeroCount.short_zeta_unit_bin_sum_le_jensen; MathCollab.Density.LocalZeroCount; MathCollab.Density.ZetaGrowth}

\begin{lemma}[Uniform logarithmic local multiplicity and extraction]\label{lem:representatives}
There is $B>0$ such that for all $\sigma\ge9/13$, $H\ge2$, every unit
bin in $Z(\sigma;H,2H)$ has total multiplicity at most $B\log(2H)$.
There is therefore a subset $R\subset Z(\sigma;H,2H)$ with injective,
one-separated ordinates and
\[
 n(\sigma;H,2H)\le2B\log(2H)|R|.
\]
\end{lemma}
\begin{proof}
For large $H$, Input~\ref{input:jensen} is at most a fixed multiple of
$\log(2H)$. For bounded $2\le H\le H_0$, every bin sum is at most the
finite count $Q=N(9/13,2H_0)$; since $\log(2H)\ge\log4$, add
$Q/\log4$ to the uniform coefficient.

For the finite weighted extraction, partition the occupied unit bins into
even and odd classes and choose the class carrying at least half the total
weight. Choose one zero from each occupied bin in that class. Different
chosen bins are separated by at least two integer bin indices; their
ordinates are at least one apart and cannot coincide. A bin has weight at
most $B\log(2H)$, so the total weighted count is at most
$2B\log(2H)$ times the number of representatives. The source packages this
as weighted separated extraction, preserving full zero multiplicity.
\end{proof}
\source{GuthMaynard.Improvement.ShortLocalZeroCount.short_zeta_local_multiplicity_bound, short_zeta_separated_representatives; MathCollab.Density.ZeroSeparation}

\begin{proposition}[Actual slab density in the two closed windows]\label{prop:slabs}
For every $\eps>0$, one constant and one eventual height threshold work
for all $\sigma$ in either specified interval:
\[
 n(\sigma;H,2H)\le C_\eps H^{g(\sigma)+\eps}
    \quad(9/13\le\sigma\le5/7),
\]
\[
 n(\sigma;H,2H)\le C_\eps H^{h(\sigma)+\eps}
    \quad(5/7\le\sigma\le3/4).
\]
\end{proposition}
\begin{proof}
Choose representatives as in Lemma~\ref{lem:representatives}.
Filter them into $R_I=\{\rho:|F_H(\rho)|\ge1/3\}$ and its complement
$R_{II}$. The detector alternative implies that every complement member
has $|I(\rho;X_H,Y_H)|\ge1/3$. Both subsets preserve separation and
ordinate injectivity. Apply Proposition~\ref{prop:typeI} with loss
$\eps/2$ and Proposition~\ref{prop:typeII} with loss $\eps/2$, using its
exact exponent comparisons. Then
\[
 |R|\le(C_I+C_{II})H^{q(\sigma)+\eps/2},
 \qquad q=g\text{ or }h.
\]
Multiplicities are restored by the preceding logarithmic factor.
Since $\log(2H)\le(\log2+2/\eps)H^{\eps/2}$, the total exponent is
$q(\sigma)+\eps$. The constant is selected before $\sigma,H$.
\end{proof}
\source{GuthMaynard.Improvement.ShortDetectorCounting.short_slab_count_of_separated_branch_bounds; WeakShortDensity.short_slab_density; UpperCurveDensity.upper_curve_slab_density}

\section{From eventual positive slabs to both-sign all-height counts}
The source proves conjugation symmetry of zeta and invariance of its analytic
multiplicity, hence
$N(\sigma,T)\le2n(\sigma;0,T)$.
The inequality, rather than equality, permits the real-axis overlap.
A closed-slab split satisfies
$n(\sigma;0,2H)\le n(\sigma;0,H)+n(\sigma;H,2H)$;
zeros on the shared boundary can be counted twice only on the upper-bound side.
\source{MathCollab.Density.ZetaConjugation}

\begin{lemma}[Exact dyadic summation with positive exponent]\label{lem:dyadic}
For fixed $\sigma$, if $p>0$ and eventually
$n(\sigma;H,2H)\le CH^p$, $C>0$, then one $A>0$ works for all $T\ge2$:
$N(\sigma,T)\le AT^p$.
\end{lemma}
\begin{proof}
First extend the slab estimate to every $H\ge1$.
Choose $R=\max(1,H_0)$ from its eventual threshold and put
$Q=n(\sigma;0,2R)$. For $1\le H<R$, monotonicity gives the slab count
at most $Q\le QH^p$; for $H\ge R$ use $CH^p$.
Thus $B=C+Q$ works on every such slab.
Put $r=2^p>1$, $Q_0=n(\sigma;0,1)$ and
$D=Q_0+B/(r-1)>0$. The source notes $D+B\le Dr$.
Induction using the closed-slab split gives
\[
 n(\sigma;0,2^j)\le D(2^j)^p.
\]
For $T\ge1$ choose $2^j\le T<2^{j+1}$. Monotonicity gives
$n(\sigma;0,T)\le D(2^{j+1})^p\le DrT^p$.
Conjugation bounds $N(\sigma,T)$ by twice that quantity. The geometric
constant $1/(2^p-1)$ is finite because $p>0$. No epsilon is spent in
this summation; its dependence on $\sigma$ is retained.
\end{proof}
\source{MathCollab.Density.DensityDyadicSummation: zetaSlab_extend_power_bound, zetaPositive_bound_of_slabs, zetaDensity_bound_of_eventual_slabs}

\begin{corollary}[Fixed-line density in the two new windows]\label{cor:windows}
For $9/13\le\sigma\le5/7$, every $\eps>0$ gives
$N(\sigma,T)\ll_{\sigma,\eps}T^{g(\sigma)+\eps}$ for all $T\ge2$.
For $5/7\le\sigma\le3/4$ the same assertion holds with $h(\sigma)$.
\end{corollary}
\begin{proof}
The exponents $g(\sigma)+\eps$ and $h(\sigma)+\eps$ are strictly
positive on these closed intervals. Apply Lemma~\ref{lem:dyadic} to
Proposition~\ref{prop:slabs} for the fixed line. The source does not infer
uniformity in $\sigma$ from this fixed-line geometric summation.
\end{proof}
\source{GuthMaynard.Improvement.WeakShortDensity.short_density; UpperCurveDensity.upper_curve_density}

\chapter{The previously established lower and strictly upper branches}\label{ch:otherbranches}
\section{The native Ingham branch and count-convention bridge}
The lower input is not taken from the informal name ``Ingham'' alone.
The development proves a positive-slab theorem through its native classical
detector, then translates the native rectangle count into the actual count
of Chapter~1.

\begin{definition}[Classical exponent certificate]
The native logarithmic scale is $\tau=\log T/\log Q$ when $Q>1$.
Put
\[
 E_{\rm MHH}(\sigma,\tau)=\max(2-2\sigma,
             \min(\tau+1-2\sigma,\tau+4-6\sigma)),
\]
\[
 E_{\rm sub}(\sigma,\tau_0,\tau)=
             \max(2-2\sigma,3-3\sigma+\tau-\tau_0).
\]
The source structure \lean{EndpointScaleCertificate} consists of
$\tau_0>0$, $3-3\sigma\le\tau_0$, and the following four obligations:
\begin{enumerate}[label=(\roman*)]
\item $E_{\rm MHH}\le2-2\sigma$ for $0\le\tau\le\tau_0+\sigma-1$;
\item $E_{\rm sub}\le(3-3\sigma)\tau/\tau_0$ for
$2\tau_0/3\le\tau\le\tau_0$;
\item for $4\tau_0/3\le\tau\le8\tau_0/3$, either $\tau/2$ or $\tau/3$
is in $[2\tau_0/3,\tau_0]$;
\item the exceptional zeta-polynomial interval
$2\le\tau<4\tau_0/3$ is empty or lies below $6\sigma-3$,
the proved Weyl exclusion threshold.
\end{enumerate}
\end{definition}
\source{GuthMaynard.NativeMoment.FiniteScaleAssembly: classicalMHHExponent, corollary1110Envelope, EndpointScaleCertificate, ZetaWindowResolved}

\begin{lemma}[The Ingham certificate]
For $1/2<\sigma<1$, $\tau_0=2-\sigma$ satisfies this certificate.
\end{lemma}
\begin{proof}
$\tau_0>0$ and $3-3\sigma\le2-\sigma$ follow from $\sigma>1/2$.
In (i) the upper bound is $1$, and the second-moment branch of
$E_{\rm MHH}$ is at most $2-2\sigma$ for $\tau\le1$.
For (ii), the first component fits the target by
$\tau\ge2\tau_0/3$. For the second, multiplying by $\tau_0>0$
reduces the comparison to
$(\tau_0-(3-3\sigma))(\tau_0-\tau)\ge0$.
For (iii), choose the half-scale if $\tau\le2\tau_0$, otherwise the
third-scale. For (iv), $4(2-\sigma)/3<2$ when $\sigma>1/2$, so the
interval is empty. These are exactly the finite exponent obligations;
they do not themselves supply the analytic detector estimate.
\end{proof}
\source{GuthMaynard.NativeMoment.FiniteScaleAssembly.ingham_endpoint_scale_certificate}

\begin{analyticinput}[Native classical endpoint positive-slab theorem]\label{input:inghamnative}
For $1/2<\sigma<1$ and an endpoint certificate at $\tau_0$, the actual
native closed-rectangle count satisfies
\[
 \forall\eps>0,\qquad
 n_{\rm nat}(\sigma;T,2T)\ll_{\sigma,\eps}T^{3(1-\sigma)/\tau_0+\eps}
 \quad\text{eventually in }T.
\]
At $\sigma=1/2$, the native global count separately obeys an
$\eps$-power bound relative to $T$.
\end{analyticinput}
\paragraph{The analytic content consumed by this input.}
The native finite detector dichotomy first gives zero count zero, a
Type-I finite polynomial witness, or a Type-II witness. The Type-I
logarithmic scale is split into the low range, the basic range
$[2\tau_0/3,\tau_0]$, the raised range $\tau\ge4\tau_0/3$, and the
remaining medium range. In the latter the source separates
$\tau\le11/10$ from $11/10<\tau<4\tau_0/3$.
The low/basic/powered consumers use actual MHH, convolution-power,
coefficient normalization, and displacement/reseparation bounds.
The last medium range uses the proved native ordinary-zeta reflection
consumer \lean{classical_medium_typeI_witness_consumer_native}; it is
not silently eliminated by the empty exceptional-zeta window.
The Type-II consumer uses the actual fourth-moment reduction.
The native source selects the maximum of the finitely many branch constants
and a common maximum of their height thresholds, retaining the binwise
analytic-multiplicity factor. This yields the positive-slab estimate.
The half-line endpoint uses the separately proved global zero-count estimate.
\source{GuthMaynard.NativeMoment.ClassicalEndpointDensity.classical_endpoint_positive_slab_native, ingham_endpoint_positive_slab_native; MediumTypeIEndpoint; TypeIIFourthMomentReduction; ClassicalDensity.ingham_zero_density_at_half_native}

\begin{proposition}[The actual Ingham lower strip]\label{prop:ingham}
For each $\sigma\in[1/2,9/13]$ and $\eps>0$ there is $C>0$ such that
for all $T\ge2$,
\[
 N(\sigma,T)\le CT^{3(1-\sigma)/(2-\sigma)+\eps}.
\]
\end{proposition}
\begin{proof}
For $\sigma>1/2$ use the certificate just proved and
Input~\ref{input:inghamnative}, followed by the native dyadic-to-global
transfer. The native bound is expressed as
\lean{EpsilonPowerBound}: for each positive $\eps$ the absolute count
is Big-O at infinity of $T^\eps|T^p|$.
The native rectangle allows $\Rea\rho\le1$ rather than $\Rea\rho<1$.
Since $\sigma>0$, every native zero has positive real part, and imported
zeta nonvanishing for $\Rea\rho\ge1$ makes its upper inequality strict.
Thus the finite native zero set equals $Z(\sigma;-T,T)$ exactly, and both
sums use the same natural analytic order. At $\sigma=1/2$ use the
separate native global estimate, where $p=1$.
For either case a positive Big-O constant is obtained; because
$p=3(1-\sigma)/(2-\sigma)\ge0$, the finite bounded-height range is
absorbed by $Q=N(\sigma,R)$ with $R\ge2$ above the eventual threshold.
$C+Q$ works since $T^{p+\eps}\ge1$ for $T\ge2$. This proves the
all-height actual-count statement without changing multiplicity conventions.
\end{proof}
\source{GuthMaynard.Improvement.NativeIngham.ingham_lower_density; NativeDensitySourceBridge.native_N_eq_zetaDensityCount, native_density_epsilon_bound; NativeDensityConventionBridge}

\section{The physical twelfth moment and the high-value producer}
Put $z(t)=|\zeta(1/2+\ii t)|$ and
$E(H,V)=\{t\in[H,2H]:z(t)\ge V\}$.

\begin{analyticinput}[Native high-value finite-peak theorem]\label{input:peaks}
For every $\eta>0$, eventually in $H$, uniformly for
$V>0$, $V\ge H^{1/8+\eta}$, every finite one-separated
$W\subset E(H,V)$ obeys
\[
 |W|V^{12}\le H^{2+\eta}.
\]
\end{analyticinput}
This actual-zeta input is supplied by the native point-mean producer and the
native Atkinson stationary-source local-mean producer:
\lean{PointMean.pointMeanInput_native} and
\lean{Atkinson.localMeanStationaryInput_native}. Their packet width ranges,
local means, and stationary transform estimates are substantive analytic
foundations. They are not deduced in this document from a pointwise growth
bound or an unnamed moment assumption.
\source{MathCollab.Density.Stronger.NativePeaks.pointValue_peak_card_native; PointMean.PeaksInput; Atkinson.AtkinsonNativeLocalMean; Peaks.WidthRanges}

\paragraph{Exact point-mean and packet interfaces.}
Write $z(t)=|\zeta(1/2+\ii t)|$. The proved point-mean interface selects
$P>0$ such that for $T\ge10$, $V>0$, $G,L\ge0$, a finite one-separated
$W\subset[c-G/2,c+G/2]$, with $10\le c-G/2$, $c+G/2\le T$,
$(\log t)^2\le L$ on $W$, $G/2+L\le G$, and $z(t)\ge V$, one has
\[
 V^2|W|\le P\log T\left(|W|+4\int_{c-G}^{c+G}z(u)^2du\right).
\]
Define the actual excess
$E_G(t;e)=\max(0,\int_{t-G}^{t+G}z(u)^2du-e)$.
The proved packet interface, for every $\delta,\kappa,\nu>0$, selects
$C,D>0,H_0\ge40000$ before all subsequent variables. If $H\ge H_0$,
$W$ is $G$-separated, $H\le t\le2H$ and
$t^\delta\le G\le t^{1/2-\delta}$, $t^{1/4+\kappa}\le G$ on $W$,
and $W\subset[A,A+L]$, it states
\[
 \left(\sum_{t\in W}E_G(t;CG\log t)\right)^2
 \le DH^\nu\left(|W|H/G+|W|^2\sqrt{GL}\right).
\]
These are the precise two imported analytic inequalities in the peak argument.
Their producers use the actual Gamma--zeta-square point contour and the
stationary Atkinson local mean, respectively.
\source{MathCollab.Density.Stronger.Peaks.Inputs: PointMeanInput, LocalMeanPacketInput; PointMean.Equation44.heathBrown_equation44_native}

\paragraph{Finite packet absorption and peak aggregation.}
Given a packet coefficient $A_0>0$ and excess level $Y_0>0$, choose
$L_0=(Y_0^2/(4A_0))^2/G$. Then
$2A_0\sqrt{GL_0}=Y_0^2/2\le Y_0^2$.
For a localized set of size $r$, $rY_0\le\sum E_G$ and the packet inequality
imply $r^2Y_0^2\le A_0(rH/G+r^2\sqrt{GL_0})$.
If $r>0$, absorption and cancellation give
$r\le2A_0H/(GY_0^2)$; if $r=0$ the same bound is immediate.
Partition $[H,2H]$ into at most $\lfloor H/L_0\rfloor+1$ fibers.
This gives the exact global bound
\[
 r\le2A_0H/(GY_0^2)+32A_0^3H^2/Y_0^6.
\]
Applying the packet interface with loss $\nu/3$ puts both terms below
$D' H^\nu(H/(GY_0^2)+H^2/Y_0^6)$.

For point peaks, use half-width bins
$n=\lfloor(t-H)/(G/2)\rfloor_+$ with center $c_n=H+nG/2$ and occupancy
$r_n$. The point-mean input gives, with $\mathcal L=\log(3H)$,
$V^2r_n\le P\mathcal L(r_n+4I_n)$, where
$I_n=\int_{c_n-G}^{c_n+G}z^2$.
When $V^2\ge2P\mathcal L$, put $A_1=V^2/(16P\mathcal L)$;
then $I_n\ge2r_nA_1$.
If the packet error is at most $A_1$, every bin with $r_n\ge m\ge1$
has excess at least $mA_1$. Even and odd bin centers separately are
$G$-separated. The preceding excess-count bound yields
\[
 |\{n:r_n\ge m\}|\le2D'H^\nu
       \left(\frac H{G(mA_1)^2}+\frac{H^2}{(mA_1)^6}\right)
 \le\frac{2D'H^\nu}{m^2}
       \left(\frac H{GA_1^2}+\frac{H^2}{A_1^6}\right).
\]
The finite occupancy identity $|W|=\sum_{m\ge1}|\{n:r_n\ge m\}|$ and
$\sum_{m\ge1}m^{-2}\le2$ give the point-count bound at width $G$.
Choose $K=16PC+1$ and $G=V^2/(K\mathcal L^2)$. The error condition follows
from $K\ge16PC$; the point-mean width condition follows eventually from
$2\mathcal L^2\le G$. Subject to
\[
 K\mathcal L^2(2H)^{1/4+\kappa}\le V^2
       \le K\mathcal L^2H^{1/2-\delta},
\]
the exact width conditions hold throughout the slab, and substitution yields
\[
 |W|\le D''H^\nu
       \left(H\mathcal L^4/V^6+H^2\mathcal L^6/V^{12}\right).
\]
For the requested peak loss $\eps$, the source chooses
$\eta=\min(\eps/16,1/1000)$, $\delta=1/48$, $\kappa=\nu=\eta$.
For nonempty $W$, its actual Weyl-growth bound gives $V\le H^{1/6+\eta}$.
The lower threshold $V\ge H^{1/8+\eps}$ and this growth bound make both
width comparisons valid eventually. Multiplying by $V^{12}$ gives
\[
 D''\mathcal L^4H^{2+7\eta}+D''\mathcal L^6H^{2+\eta}
 \le H^{2+8\eta}\le H^{2+\eps}
\]
eventually. Empty sets are handled directly. This proves the finite peak
consumer used in Input~\ref{input:peaks}, including every occupancy and
width loss, relative to its two precisely stated analytic producers.
\source{MathCollab.Density.Stronger.AtkinsonCardinalityAbsorption; AtkinsonGlobalCardinality; Peaks.Clusters; Peaks.ClusterCount; Peaks.WidthRanges}

\begin{lemma}[High-value measure and the twelfth moment]\label{lem:twelfth}
For each $\eps>0$ the source proves
\[
 \int_H^{2H}z(t)^{12}dt\ll_\eps H^{2+\eps}
 \quad\text{eventually},\qquad
 \int_0^{3T}z(t)^{12}dt\ll_\eps T^{2+\eps}
 \quad\text{eventually}.
\]
\end{lemma}
\begin{proof}
The finite separated-cover lemma applied to Input~\ref{input:peaks}
gives $|E(H,V)|\le2H^{2+\eta}/V^{12}$ for $V\ge H^{1/8+\eta}$.
For a requested exponent $\eps$ take $\eta=\eps/20$.
Below $V=H^{1/8+\eta}$, $z^{12}\le V^8z^4$, so
Input~\ref{input:fourthzeta} contributes at most
$CH^{2+9\eta}$.
For the high values the source uses its established actual growth bound
$z^{12}\le H^3$ on the dyadic interval eventually. With
$\eta'<\eta$, the measure-tail bound at $\eta'$ and layer integration
of the twelfth power give
\[
 \int_{E(H,H^{1/8+\eta'})}z^{12}
 \le2H^{2+\eta'}\left(1+\log\frac{H^3}{(H^{1/8+\eta'})^{12}}\right)
 \le2H^{2+\eta'}(1+3\log H).
\]
The logarithm is absorbed into the positive exponent gap to $2+\eta$.
The set with the larger cutoff $H^{1/8+\eta}$ is a subset, so its
high contribution is at most $H^{2+\eta}$ eventually.
Since $9\eta\le\eps$, summing the two parts proves the dyadic bound.
The physical interval follows from the proved dyadic integral accumulation
for the continuous nonnegative function $z^{12}$, with exponent
$2+\eps>0$, absorbing the initial compact interval. Replacing the terminal
height by $3T$ changes only the constant by $3^{2+\eps}$.
\end{proof}
\source{MathCollab.Density.Stronger.ConditionalMoment.zeta_twelfth_dyadic_of_fourth_and_tail, zeta_twelfth_physical_of_dyadic; MomentGrowth; NativeMoment.zeta_twelfth_physical_native}

\section{The smooth detector family above three quarters}
For a fixed $3/4<\sigma<1$ put
\[
 \alpha=1/(2\sigma),\quad c=3/(4\sigma)=3\alpha/2,\quad
 q(\sigma)=3(1-\sigma)/(2\sigma),\quad X=T^\delta,\quad Y=T^c.
\]
We have $0<c<1$, and choose $\delta>0$ before heights and zeros.
The source's actual smooth dyadic partition is constructed from mathlib's
smooth transition $s$, equal to zero on $(-\infty,0]$ and one on $[1,\infty)$:
\[
 d(x)=1-s(4x-5),\qquad \omega(x)=d(x)-d(2x).
\]
It is nonnegative, at most one, and supported in $[5/8,3/2]\subset(1/2,2)$.
\begin{lemma}[The exact signed-scale partition]\label{lem:signedpartition}
For $x>0$,
\[
 \sum_{j=0}^{\infty}\omega(x/2^j)=1-d(2x),\qquad
 \sum_{n=0}^{\infty}\omega(x2^{n+1})=d(2x),\qquad
 \sum_{j\in\Z}\omega(x/2^j)=1.
\]
The second sum is the sum over the negative integer scales $j=-(n+1)$;
it does not include the scale $j=0$.
\end{lemma}
\begin{proof}
By the definition $\omega(y)=d(y)-d(2y)$, for every natural $m$ the
nonnegative-scale partial sum telescopes exactly to
\[
 \sum_{j=0}^{m}\omega(x/2^j)=d(x/2^m)-d(2x).
\]
As $m\to\infty$, $x/2^m\to0$ and $d(x/2^m)\to d(0)=1$.
For the negative scales, substitute $j=-(n+1)$, so that
$x/2^j=x2^{n+1}$. Their partial sum is exactly
\[
 \sum_{n=0}^{m}\omega(x2^{n+1})
 =\sum_{n=0}^{m}\bigl(d(x2^{n+1})-d(x2^{n+2})\bigr)
 =d(2x)-d(x2^{m+2}).
\]
Since $x>0$, $x2^{m+2}\to\infty$; eventually it is at least $3/2$, where
$d$ is zero. Thus this limit is $d(2x)$, not $d(x)$.
All summands are nonnegative. Combining the disjoint index sets
$\{0,1,\ldots\}$ and $\{-1,-2,\ldots\}$ gives
$[1-d(2x)]+d(2x)=1$, with neither a missing nor a duplicated zero scale.
For fixed $x>0$ the support condition also bounds all nonzero summands
between finite integer limits, so the integer-scale sum is pointwise finite.
\end{proof}
\source{MathCollab.Density.Stronger.SmoothPartition: hasSum_smoothDyadicWeight_nonneg_scales, hasSum_smoothDyadicWeight_neg_scales, hasSum_smoothDyadicWeight, smoothDyadicWeight_pointwise_finite}
Splitting a block at its half-scale endpoints introduces no endpoint
contribution because the support is strictly contained in $(1/2,2)$.
Its actual block is
\[
 B_j(\rho)=\sum_{n\ge2}c_X(n)n^{-\rho}\e^{-n/Y}\omega(n/2^j).
\]
The finite scale set is $\mathcal S_R=\{j\in\Z:j\ge0,2^j\le R\}$,
$R=Y(\log T)^2$. A long block has $2^j>T^{c/2}$ and
$|B_j(\rho)|\ge1/(24\log T)$.
For the short family split a smooth block into natural dyadic half scales
$Q_{j,e}=2^{j_{\rm nat}-1+e}$ for $e=0,1$.
The coefficient on a half is
$c_X(n)\e^{-n/Y}\omega(n/2^j)$, retaining the original smooth scale $2^j$.
Apply the same real-part Taylor weights of Chapter~\ref{ch:detector} to
each half. The Taylor index set is
$\mathcal S_R\times\{0,1\}\times\{0,\ldots,\lceil10\log T\rceil_+\}$.
Let $K=\lceil2\alpha/\delta\rceil_+$.
Take convolution powers $1\le k\le K$ and pieces $0\le\ell<k$,
retaining only members with $T^\alpha\le Q^k\le T^c$.
Their integer starts are $L_q=2^\ell Q^k$, their coefficients are the
conjugated convolution coefficients divided by $T^\eta$.
This exactly describes \lean{smoothPoweredIndices} and
\lean{smoothPoweredFamilyCoefficient}.
\source{MathCollab.Density.Stronger.SmoothDetector; ShortFixedFamily; ShortTaylorCover; ShortPoweredFamily}

\begin{analyticinput}[Native smooth detector cover and long-mean reduction]\label{input:smoothcover}
For fixed $3/4<\sigma<1$, $0<\delta\le1/8$, $\eta>0$, eventually every
actual zero in the positive slab is either covered by the preceding
normalized powered short family at its original ordinate, with value
$\ge L_q^\sigma T^{-2\eta}$, or has a long smooth block.
The complementary, multiplicity-weighted long count $N_{\rm long}$ satisfies
\begin{equation}\label{eq:longraw}
 N_{\rm long}Y^{6\sigma-3}
 \le CT^{6\delta}(\log T)^{13}
          \left(\int_0^{3T}z(t)^{12}dt+T^{-113}\right).
\end{equation}
\end{analyticinput}
\paragraph{What is used in the long reduction.}
For the complementary zero the source proves the lower bound
\[
 Y^{(\sigma-1/2)/2}\le48C\sqrt X\log T\,
       \int_\R\frac{z(\Ima\rho+t)}{(1+|t|)^{128}}\,dt.
\]
It extracts separated representatives with local analytic multiplicity,
raises this bound to the twelfth power, and applies H\"older and the
physical weighted-mean packing lemma. The near weighted integral is
controlled by the physical twelfth moment; the far tail by the proved zeta
growth and the integrable weight of order $128$. This yields the exact
exponents $X^6$, $Y^{6\sigma-3}$, $(\log T)^{13}$, and remainder
$T^{-113}$ in \eqref{eq:longraw}. Long and short counts form an exact
partition of the original slab, weighted by $m(\rho)$.
\source{MathCollab.Density.Stronger.ShortPoweredFamily.smooth_long_or_short_powered_cover; LongMultiplicityReduction.smoothLongZeroCount_physical_moment_reduction; WeightedCriticalMean; LongMeanCounting}

\begin{proposition}[Short part of the stronger density branch]\label{prop:strongshort}
If $\eta>0$ and $8\eta<\sigma-3/4$, the actual multiplicity count
covered by the smooth powered family satisfies
$N_{\rm short}\ll_{\sigma,\delta,\eta}T^{q(\sigma)+12\eta}$ eventually.
\end{proposition}
\begin{proof}
The source establishes for every eligible member
$T^\alpha\le L_q\le2^KT^c\le T$ eventually, bounded normalized
coefficients, and at most $144K^2(\log T)^2$ members.
Use the earlier native integer large-values bound (the high-threshold
input described after \eqref{eq:classical}) at comparison height
$H_a=32T(\log T)^2$. It requires $H_a^{1/3}\le L_q\le H_a$ and
$L_q^{3/4+\kappa}\le L_q^\sigma T^{-2\eta}$.
Take $\kappa=(\sigma-3/4)/2>0$; the margin $8\eta<\sigma-3/4$ implies
$3/4+\kappa\le\sigma-4\eta$. Since $L_q\ge T^\alpha\ge\sqrt T$,
these height and threshold conditions hold eventually exactly as in the
source.

With $V_q=L_q^\sigma T^{-2\eta}$ the two terms are
$L_q^{2-2\sigma}T^{4\eta}$ and
$L_q^{3-4\sigma}\sqrt{H_a}\,T^{8\eta}$.
The positive length exponent uses the upper length bound and
$2c(1-\sigma)=q(\sigma)$; the negative one uses the lower bound and
$1/2+(3-4\sigma)\alpha=q(\sigma)$.
Each member's budget is at most
$(2^K+32)\log T\,T^{q(\sigma)+8\eta}$.
The indexed finite cover and multiplicity extraction add a factor
$\log(2T)H_a^\eta$ and the number of family members.
Eventually $H_a^\eta\le T^{2\eta}$, and the resulting four logarithms
are at most $T^{2\eta}$. This gives exactly the total loss $12\eta$.
It bounds only the explicitly covered subset; the other zeros remain in
$N_{\rm long}$.
\end{proof}
\source{MathCollab.Density.Stronger.ShortExponentAccounting.smooth_family_sum_eventually; ShortZeroCount.smoothShortZeroCount_bound}

\begin{proposition}[The strictly upper density branch]\label{prop:strong}
For every fixed $\sigma>3/4$ and $\eps>0$ there is $C>0$ such that
for all $T\ge2$,
\[
 N(\sigma,T)\le CT^{3(1-\sigma)/(2\sigma)+\eps}.
\]
\end{proposition}
\begin{proof}
For $\sigma\ge1$ the count is zero, so take $C=1$.
For $3/4<\sigma<1$, put $d=1-c>0$. Choose
\[
 \delta=\min(1/16,d/24),\quad
 \eta=\min(\eps/24,(\sigma-3/4)/16),\quad\nu=d/16.
\]
Then $0<\delta\le1/8$, $\delta<d/12$, $8\eta<\sigma-3/4$,
$12\eta\le\eps$, and $0<\nu\le d/8$.
Lemma~\ref{lem:twelfth} gives $\int_0^{3T}z^{12}\le DT^{2+\nu}$
eventually. Insert this in \eqref{eq:longraw}; its extra remainder is
also at most $T^{2+\nu}$ for $T\ge1$.
The identity
\[
 2+(6-12\sigma)c/2=q(\sigma)-(1-c)
\]
shows that the raw long exponent is $q-d+6\delta+\nu$.
Because $6\delta<d/2$ and $\nu\le d/8$, a positive margin remains to
$q-d/4$; the thirteen logarithms are absorbed into that margin.
Thus $N_{\rm long}\ll T^{q-d/4}$ eventually.
Proposition~\ref{prop:strongshort} gives
$N_{\rm short}\ll T^{q+12\eta}$. Their exact partition gives an
actual positive-slab bound with exponent $q+12\eta>0$.
Apply Lemma~\ref{lem:dyadic} and finally monotonicity in the exponent to
replace $12\eta$ by $\eps$ for $T\ge2$.
The choice $d>0$ and the high-threshold margin fail at $\sigma=3/4$;
no endpoint at that line is inferred.
\end{proof}
\source{MathCollab.Density.Stronger.FinalParameters.exists_stronger_loss_parameters; ConditionalDensity.stronger_density_bound_of_twelfth_moment; NativeDensity.stronger_density_bound_native}

\section{The separate coefficient-two branch used in the global proof}
The source's global-uniform theorem uses an earlier coefficient-two result
above $3/4$, not the stronger branch just proved. To preserve that actual
proof route we record its exact established interface.
\begin{analyticinput}[Native Weyl-based density conclusion]\label{input:two}
The source proves, for $\sigma>3/4$ and every $\eps>0$,
\[
 N(\sigma,T)\ll_{\sigma,\eps}T^{2(1-\sigma)+\eps}\quad(T\ge2).
\]
It discharges the actual critical-line Weyl input
$z(t)\le C_w(1+|t|)^{1/6+\eta_0}$ with $0\le\eta_0<1/48$.
\end{analyticinput}
\paragraph{Use of the imported machinery.}
The source chooses the positive counting loss
$\min((\sigma-3/4)/16,\eps/24)$ before heights, obtains an eventual
actual positive-slab bound with exponent $2(1-\sigma)$ plus twelve such
losses, and uses the actual-count dyadic summation. The Weyl growth is
proved by the native van der Corput, second-derivative, and explicit Weyl
chain. The contour errors, powered family, separated counts, and analytic
multiplicity restoration in that branch are proved dependencies, not
unproved density premises. If $\sigma\ge1$ it again uses the empty count.
\source{MathCollab.Density.DensityConclusion.zeta_density_bound_of_weyl; DensityTheorem.zeta_density_bound; GuthMaynard.WeylZeta; WeylPort}

\chapter{Endpoint assembly, uniform constants, and the boundary audit}\label{ch:final}
\section{The fixed-line piecewise theorem}
\begin{proof}[Proof of Theorem~\ref{thm:full}]
If $1/2\le\sigma\le9/13$, apply Proposition~\ref{prop:ingham}.
If $9/13<\sigma\le5/7$, use Corollary~\ref{cor:windows} with $g$.
If $5/7<\sigma\le3/4$, use its $h$ bound.
If $3/4<\sigma\le1$, use Proposition~\ref{prop:strong}.
The resulting exponent is exactly $f(\sigma)=A(\sigma)(1-\sigma)$.
The constant in each branch may depend on the fixed $\sigma$ and $\eps$;
this case distinction makes no additional uniformity claim.
\end{proof}
\source{GuthMaynard.Improvement.CombinedDensityCurve.combined_density_curve; GuthMaynardStatement.coefficient_times_gap, full_curve}

The endpoint values are
\[
 f(1/2)=1,\quad f(9/13)=12/17,\quad f(5/7)=9/14,
 \quad f(3/4)=9/16,\quad f(1)=0.
\]
At $9/13$ the two rational formulas coincide; at $5/7$ the rational
formula equals $(9/4)(1-\sigma)$. The value at $3/4$ belongs to the
closed upper window, not the strictly upper branch.
\source{GuthMaynard.Improvement.CombinedDensityCurve.combined_density_curve_endpoints}

\section{The actual route to the fixed-line coefficient 39/17}
Let $c_0=39/17$. The source compares
\[
 \frac{3(1-\sigma)}{2-\sigma}\le c_0(1-\sigma)
       \quad(1/2\le\sigma\le9/13),
\]
\[
 \frac{18(1-\sigma)}{7\sigma+3}\le c_0(1-\sigma)
       \quad(9/13\le\sigma\le5/7).
\]
The first uses $2-\sigma\ge17/13$, and the second
$7\sigma+3\ge102/13$. Both denominators are positive and $1-\sigma\ge0$.
Above $3/4$ Input~\ref{input:two} fits since $2<c_0$ and
$1-\sigma\ge0$.
The source closes $[5/7,3/4]$ through its \emph{separate} padded upper-window
route. It uses the same fixed powered detector family at lower exponent
$3/4$, upper exponent $9/8$, but uses the fixed second split $7/8$ rather
than $s(\sigma)$. The needed comparisons are
\[
 \tfrac94(1-\sigma),\quad39/14-3\sigma,\quad
 63/20-7\sigma/2,\quad15/8-7\sigma/4
       \ \le c_0(1-\sigma).
\]
They are affine in $\sigma$ and hold at both endpoints $5/7$ and $3/4$.
Padded time $S=H$ below $H^{5/6}$ and $S=N^{6/5}$ between
$H^{5/6}$ and $H^{7/8}$ yields the first three budgets; the mean-value
branch above $H^{7/8}$ yields the fourth. The loss is $10\delta$ as in
Chapter~\ref{ch:windows}. The fixed-family cover adds $\eta$, the
unconditional weak Type-II exponent fits $c_0(1-\sigma)$, and Jensen
restoration plus dyadic summation gives the actual upper density bound.
Thus the whole fixed-line $c_0$ theorem has no
\lean{GuthMaynardUpperDensity} or \lean{ShortMobiusFourthMoment} premise.
\source{GuthMaynard.Improvement.UpperWindowBudgets; UpperWindowLargeValues; UpperTypeICount; WeakUpperDensity.upper_density; GlobalDensityIngham.global_density_fixed_sigma_of_ingham}

\section{Finite-grid uniformization}
\begin{lemma}[One-sided-modulus uniformization]\label{lem:grid}
Let $a<b$, $K>0$, and $q:[a,b]\to\R$ satisfy
\[
 q(x)\le q(y)+K(y-x)\quad(a\le x\le y\le b).
\]
Suppose each $\sigma\in[a,b]$ and each $\eta>0$ has a constant with
$N(\sigma,T)\le C_{\sigma,\eta}T^{q(\sigma)+\eta}$ for all $T\ge2$.
For every $\eps>0$ a single positive $C$ then works for all $\sigma,T$:
$N(\sigma,T)\le CT^{q(\sigma)+\eps}$.
\end{lemma}
\begin{proof}
Choose a natural $n>2K(b-a)/\eps$. Set
$x_j=a+j(b-a)/n$ for $j=0,\ldots,n$.
At each grid point choose $C_j>0$ with the fixed-line estimate at loss
$\eps/2$. Let $C=1+\sum_{j=0}^n C_j$.
For a given $\sigma$, let
$j=\lfloor(\sigma-a)n/(b-a)\rfloor_+$; then $0\le j\le n$,
$x_j\le\sigma$ and $\sigma-x_j\le(b-a)/n$.
The one-sided modulus gives
$q(x_j)+\eps/2\le q(\sigma)+\eps$.
Monotonicity of the actual zero count in its real-part cutoff and
monotonicity of $T$-powers in their exponent, for $T\ge2$, give
\[
 N(\sigma,T)\le N(x_j,T)
 \le C_jT^{q(x_j)+\eps/2}
 \le C_jT^{q(\sigma)+\eps}\le CT^{q(\sigma)+\eps}.
\]
The constant was selected from a finite grid before $\sigma,T$.
\end{proof}
\source{GuthMaynard.Improvement.CurveUniformization.density_curve_uniform_on_interval}

\begin{proof}[Proof of Theorem~\ref{thm:global}]
Use the actual fixed-line coefficient-$c_0$ proof from the preceding
section. The source specializes finite-grid uniformization to
$q(\sigma)=c_0(1-\sigma)$ on $[1/2,1]$.
More explicitly choose $n>c_0/\eps$ and
$x_j=1/2+j/(2n)$, using fixed-line loss $\eps/2$.
Then $c_0(\sigma-x_j)\le c_0/(2n)\le\eps/2$.
The constant $1+\sum C_j$ and the exact comparison in Lemma~\ref{lem:grid}
give the assertion, including $\sigma=1$.
This proof does not attempt to uniformize the discontinuous exact curve.
\end{proof}
\source{GuthMaynard.Improvement.ShortDensityUniform.zetaDensityCount_uniform_of_fixed_sigma; GlobalDensityIngham.global_density_uniform_of_ingham; GlobalDensity.global_density_uniform; GuthMaynardStatement.global_uniform}

\section{Uniformity on the exact left curve and away from the strict join}
The source proves each smooth branch has a one-sided modulus with $K=4$.
For $1/2\le x\le y\le9/13$,
\[
 \frac{3(1-x)}{2-x}-\frac{3(1-y)}{2-y}
       =\frac{3(y-x)}{(2-x)(2-y)}\le4(y-x).
\]
For $9/13\le x\le y\le5/7$,
\[
 \frac{18(1-x)}{7x+3}-\frac{18(1-y)}{7y+3}
       =\frac{180(y-x)}{(7x+3)(7y+3)}\le4(y-x),
\]
using $(7x+3)(7y+3)\ge49$.
For the linear branch the difference is $(9/4)(y-x)\le4(y-x)$.
For $3/4\le x\le y\le1$ the upper rational expression has
\[
 \frac{3(1-x)}{2x}-\frac{3(1-y)}{2y}
       =\frac{3(y-x)}{2xy}\le4(y-x),\qquad xy\ge9/16.
\]
Only this modulus statement includes $x=3/4$; the stronger density theorem
itself still requires strict inequality.
\source{GuthMaynard.Improvement.CurveModulus}

\begin{proposition}[Uniform exact left curve]\label{prop:leftuniform}
For every $\eps>0$ one constant works for all $\sigma\in[1/2,3/4]$
and $T\ge2$ in the bound $N(\sigma,T)\le CT^{f(\sigma)+\eps}$.
\end{proposition}
\begin{proof}
Apply Lemma~\ref{lem:grid} separately to the three closed intervals
$[1/2,9/13]$, $[9/13,5/7]$, and $[5/7,3/4]$, using the branch
fixed-line bounds and the preceding moduli. The sum of the three positive
constants works everywhere on their union. The definitions agree at the
two interior joins, so there is no endpoint discrepancy.
\end{proof}
\source{GuthMaynard.Improvement.FullCurveLeftUniform.full_density_curve_uniform_left}

\begin{proof}[Proof of Theorem~\ref{thm:away}]
Let $C_1$ be the left constant from Proposition~\ref{prop:leftuniform}.
If $3/4+\delta<1$, apply Lemma~\ref{lem:grid} to
$[3/4+\delta,1]$ with $q(\sigma)=3(1-\sigma)/(2\sigma)$ and $K=4$.
Proposition~\ref{prop:strong} applies at every point of that interval,
so a right constant $C_2$ exists. Take $C_1+C_2$ and choose the appropriate
branch. If $3/4+\delta\ge1$, every admitted point in the right alternative
has $\sigma\ge1$, so its count is zero. The left constant alone works.
This handles intervals that degenerate to the single endpoint or are empty.
\end{proof}
\source{GuthMaynard.Improvement.CombinedDensityCurve.combined_density_curve_uniform_away_boundary; GuthMaynardStatement.uniform_away_boundary}

\section{A boundary obstruction is not a counterexample to uniform density}
\begin{proposition}[Exact boundary-modulus obstruction]
For every real $K$ there is $\sigma\in(3/4,1]$ such that
\[
 f(\sigma)+K(\sigma-3/4)<f(3/4).
\]
\end{proposition}
\begin{proof}
For every $\sigma>3/4$, the right formula satisfies
$f(\sigma)=3(1-\sigma)/(2\sigma)<1/2$, by multiplying by $2\sigma>0$.
Set $\delta_K=1/(64(|K|+1))$ and $\sigma=3/4+\delta_K$.
Then $0<\delta_K\le1/64$, $\sigma\le1$, and $K\delta_K\le1/64$.
Thus $f(\sigma)+K\delta_K<1/2+1/64=33/64<9/16=f(3/4)$.
\end{proof}
\source{GuthMaynard.Improvement.CurveBoundaryAudit.combined_curve_right_lt_half, combined_curve_boundary_modulus_obstruction}
It follows that the one-sided finite-$K$ modulus required by
Lemma~\ref{lem:grid} cannot hold across $3/4$ for the combined curve.
This proves a limitation of that uniformization method, not the falsehood
of a whole-strip exact-curve uniform density estimate.

\section{The independent actual-zero assertion in the interface}
The endpoint specification states $\operatorname{actualZeroBound}(\sigma,T,B)$
as existence of a finite set $Z$ whose membership is exactly
\[
 \zeta(\rho)=0,\quad0<\Rea\rho<1,\quad\sigma\le\Rea\rho,
 \quad|\Ima\rho|\le T,
\]
and for which $\sum_{\rho\in Z}m(\rho)\le B$.
Given any bound on $N(\sigma,T)$, choose
$Z=Z(\sigma;-T,T)$. The finite-set theorem from Chapter~1 proves its
existence, and $-T\le\Ima\rho\le T$ is equivalent to
$|\Ima\rho|\le T$. The sum defining $N$ is exactly the required sum,
so each endpoint count bound yields the independent assertion.
The coefficient identity $A(\sigma)(1-\sigma)=f(\sigma)$ is verified
branch by branch. The large-values specification is definitionally the
same bounded polynomial and the same middle-range proposition proved
in Chapter~\ref{ch:middle}. These are all the final wrapper steps.
\source{GuthMaynardStatement.actualZeroBound_of_count, coefficient_times_gap, middle_large_values in lean/Solution.lean}

\appendix
\chapter{Dependency ledger, provenance, and verification boundary}
\section{The mathematical inputs used, not additional endpoint assumptions}
The inputs below are established theorems in the pinned source closure.
Their use is explicit in the preceding proofs. A future paper expanding
all of their internal analytic proofs would require additional chapters;
this blueprint does not claim such an expansion has already been made.
\begin{enumerate}[leftmargin=*]
\item \textbf{Analytic zero foundations.} mathlib analytic zeta away from $1$,
nonvanishing in $\Rea s\ge1$, isolated zeros on compact sets, the identity
theorem, finite analytic order and local factorization. Used to make the
zero count nonvacuous and multiplicities genuine.
\item \textbf{Ordinary Fourier foundations.} Lebesgue/Bochner integration,
Tonelli/Fubini under the stated integrability, normalized Plancherel,
Fourier inversion for the cut signals, convolution regularity, rapid decay
of compact smooth/Schwartz inputs, Poisson summation with explicit scale.
Used by affine energies, source sampling, the masked cubic, and extraction.
\item \textbf{Affine native descent.} Source Lemma-9.2 recurrence, divisor
ray counting and retained Poisson budgets, adaptive smoothing, and finite
natural-number descent. The height profile, support-depth budget, and
unequal $RK$ factor are retained in Chapter~\ref{ch:affine}.
\item \textbf{Nuclear operator foundations.} Positive rank-one sums and
integrals, finite-dimensional eigenbases, witness/Bessel bounds, Schur
operator norm, cubic and mixed trace limits through finite compressions.
No infinite eigenbasis or commuting correction operator is assumed.
\item \textbf{Native classical large values.} The high-threshold integer
estimate is applied only with its $\kappa>0$ and scale restrictions; the
separated mean square closes low thresholds and real-length rounding.
\item \textbf{Actual contour foundations.} M\"obius cancellation, Mellin
inversion, regularized zeta poles, rectangle residues, Gamma strip decay,
and uniform ordinate-tail packing. Used with independent $X,Y$.
\item \textbf{Native fourth moment.} \lean{Fourth.zeta_fourth_dyadic}
feeds the weak mollified fourth moment and the high/low twelfth-moment
assembly. The sharper \lean{ShortMobiusFourthMoment} contract remains
separate and is not used in the unconditional endpoints.
\item \textbf{Native high-value peaks.} Point-mean and stationary-source
Atkinson producers plus packet width ranges feed the actual finite peak
bound. That bound feeds the high twelfth-moment tail, not an assumed
unrestricted twelfth-moment conjecture.
\item \textbf{Native Ingham endpoint.} Classical detector dichotomy and the
separate medium ordinary-zeta Type-I reflection consumer, coefficient
normalization, bounded shifts with reseparation, Type-II fourth moment,
and exact native-count bridge. The half endpoint has its own proof.
\item \textbf{Earlier coefficient-two upper density.} Explicit native Weyl
input, detector error estimates, finite powered families, multiplicity
restoration, and dyadic accumulation. It is the actual upper branch used
by \lean{GlobalDensityIngham}, even though the combined fixed-line curve
has a stronger strictly upper branch.
\end{enumerate}

\section{Proof route and compatibility names}
The final registered path is
\begin{quote}\small
\lean{GuthMaynardStatement.full_curve} follows from
\lean{combined_density_curve}.\\
\lean{GuthMaynardStatement.global_uniform} follows from
\lean{global_density_uniform}.\\
\lean{GuthMaynardStatement.uniform_away_boundary} follows from
\lean{combined_density_curve_uniform_away_boundary}.\\
\lean{GuthMaynardStatement.middle_large_values} follows from
\lean{middle_range_large_values}.
\end{quote}
The wrappers are in the file \lean{Solution.lean}; its endpoint namespace
is \lean{GuthMaynardStatement}. Similarly ``candidate'', ``conditional'', and ``gate'' in
an internal filename are not evidence of an unproved endpoint. Each explicit
intermediate contract must be checked where its producer is supplied.
The internal \lean{GuthMaynardUpperDensity} proposition is not a premise of
\lean{global_density_uniform}; \lean{InghamLowerDensity} is discharged by
\lean{ingham_lower_density}; the final short branch uses
\lean{WeakShortDensity.short_density}.

\section{Preserved attribution}
New material in the repository carries AGPL-3.0-only licensing; imported ports
retain their original notices. The pinned record identifies adapted material
from Scott McColm at revisions
\lean{6e2d10c6c2252ee1575bb7ef12dea12e2f1a3af1} (MIT-0) and
\lean{2ace9e7c09a69fdcd1edae1ab6deb7cb3b4df1be} (MIT), selected Fourier,
Poisson, Mellin and Gamma developments from Conor Grogan's
\lean{prime-minor-arcs-2-15} at
\lean{f369f267b4dcfebf010c8e7baa2c9602e2960eba} (Apache-2.0), and analytic
foundations from \lean{PrimeNumberTheoremAnd} at
\lean{4ecb950126c4290293c5662dfe0e884123171df5} (Apache-2.0).
The original repository's \lean{THIRD_PARTY_NOTICES.md}, \lean{licenses/},
source headers, and \lean{third_party/} manifests are preserved in the
archive. The informal reconstruction is not attributed to those authors as
an approved text.

\section{What has and has not been checked here}
The downloaded statement and solution files match the registry's SHA-256
values respectively:
\begin{center}
\lean{b0feb005a241a3d2da02def8c8ae83cb9db2e743b6c5679245a5c22775e7675a}\\
\lean{8b3e9e2dea3baa6c9db288f6d6bc64662b8fd1f14525599e99abb7fd9bc9f7cb}.
\end{center}
The source was read and the displayed mathematics was reconstructed and
cross-checked against the cited declarations. The LaTeX build and references
are checked separately in the accompanying validation record. This does not
constitute human peer review, an author endorsement, a novelty certificate,
or a fresh complete kernel replay of the pinned Lean development.
The registry's mechanical certificate and the repository's earlier local
verification receipt have distinct provenance; neither is silently replaced
by a claim that this reconstruction performed a full rebuild.

\section{Source references}
\begin{thebibliography}{9}
\bibitem{registry}
Palomar Registry, \emph{Zero-density estimates from an improved large-values
bound}, entry PALOMAR-2026-10-11-000010, version 1.
\url{https://palomar-registry.org/entry?id=PALOMAR-2026-10-11-000010&version=1}.
\bibitem{source}
M. Lewko, \emph{zeta-density-39-over-17-lean}, fixed repository snapshot.
\url{https://github.com/mrklwk/zeta-density-39-over-17-lean/tree/ee2a8a32be49a4f73c7386e9dffdea24d7e3d8c0}.
\bibitem{mathlib}
The mathlib community, \emph{mathlib4}, revision
\lean{065356127b1dc0016f66b7283ce0ce2c4055aa55}.
\url{https://github.com/leanprover-community/mathlib4/tree/065356127b1dc0016f66b7283ce0ce2c4055aa55}.
\bibitem{notices}
The pinned source, \emph{THIRD\_PARTY\_NOTICES.md} and source-identity
manifests, with the licenses retained for each port.
\end{thebibliography}
\end{document}
